\documentclass[a4paper,oneside,12pt]{article}
\def\runninghead{\smc }

\usepackage[russian]{babel}
\usepackage[T2A]{fontenc}
\usepackage[utf8]{inputenc}
\usepackage{amsfonts, amssymb, amsmath, amsthm}
\usepackage{multirow}
\usepackage[unicode]{hyperref}
\usepackage{ccfonts,eulervm,euler}
\renewcommand{\bfdefault}{sbc}

\theoremstyle{plain} { \swapnumbers
\newtheorem{theorem}{Теорема}[subsection]
\newtheorem{corollary}[theorem]{Следствие}
\newtheorem{lemma}[theorem]{Лемма}
\newtheorem{proposition}[theorem]{Утверждение}
\newtheorem{claim}[theorem]{Факт}
}

\theoremstyle{definition} { \swapnumbers
\newtheorem{definition}[theorem]{Определение}
\newtheorem{notation}[theorem]{Обозначение}
\newtheorem{remark}[theorem]{Замечание}
\newtheorem{remarks}[theorem]{Замечания}
\newtheorem{example}[theorem]{Пример}
\newtheorem{examples}[theorem]{Примеры}
\newtheorem{sssection}[theorem]{}
}
\pagestyle{plain}

%\newcommand{\proof}{{\par\noindent\textbf{Доказательство. }}}

%\textwidth=5cm
\oddsidemargin=-5mm
%\marginparwidth=10pt
\evensidemargin 0mm
\marginparwidth 5mm
\topmargin 0mm
\textheight 228mm
\textwidth 170mm
\headheight 0mm
\headsep 0mm
\footskip 10mm

\def\trleq{\trianglelefteq}
\def\bar{\overline}
\def\tilde{\widetilde}
\def\map{\longrightarrow}

\def\eps{\varepsilon}
\def\ph{\varphi}
\def\id{{\mathrm{id}}}
\def\GL{\operatorname{GL}}
\def\Lin{\operatorname{Lin}}
\def\Ker{\operatorname{Ker}}
\def\Image{\operatorname{Im}}
\def\End{\operatorname{End}}
\def\Aut{\operatorname{Aut}}
\def\Res{\mathrm{Res}}
\def\Cor{\mathrm{Cor}}
\def\Inf{\mathrm{Inf}}
\def\Hom{\mathrm{Hom}}
\def\la{\langle}
\def\ra{\rangle}
\def\rad{\operatorname{rad}}
\def\rk{\operatorname{rk}}
\def\ind{\operatorname{ind}}
\def\a{\alpha}
\def\b{\beta}
\def\lla{\la\!\la}
\def\rra{\ra\!\ra}

\def\dynkin#1#2#3#4#5#6#7#8{\vcenter{\vbox{\vfill
\hbox{$#1#3#4#5#6#7#8$}\nointerlineskip\vskip 3pt
\hbox{$\phantom{#1}\phantom{#3}#2$\hfil}\vfill}}}

\def\quaternion#1#2#3{\binom{{#1}\:\:{#2}}{#3}}


\begin{document}

\frenchspacing
\righthyphenmin=2

% 22.02.2010

\author{Александр Лузгарев}
\title{Алгебраическая теория квадратичных форм\thanks{Конспект лекций спецкурса, весна 2010.}}
\date{}

\maketitle

\tableofcontents
\newpage

Источники:
\par\noindent $\bullet$ Albrecht Pfister, {\it Quadratic forms with applications to algebraic geometry and topology},
London Math. Soc. Lect. Notes 217, Cambridge University Press, 1995.
\par\noindent $\bullet$ Bruno Kahn, {\it Formes quadratiques sur un corps},
Societe Mathematique de France, 2009.
\par\noindent $\bullet$ конспект лекций Олега Ижболдина, 1997.
\par\noindent $\bullet$ Philippe Gille, Tam\'as Szamuely, {\it Central simple algebras and Galois cohomology}.
Cam\-bridge University Press, 2006.

\bigskip
Автор благодарен Алексею Степанову за исправленные неточности и доказательство следствия~\ref{cor:isometry-extension}.

%%========================================================================
\section{Квадратичные формы: начало}

\subsection{Основные понятия}

\begin{definition}
Пусть $V$~--- $n$-мерное векторное пространство над полем $F$. Мы всегда
будем предполагать, что характеристика $F$ отлична от двух.
{\bf Симметричная билинейная форма} на $V$~--- это отображение
$b\colon V\times V\to k$ такое, что $b(u,v)=b(v,u)$ и $b(\a u_1+u_2,v)=\a b(u_1,v)+b(u_2,v)$.
Если $(e_1,\dots,e_n)$~--- базис $V$, то
$b(x_1e_1+\dots+x_ne_n,y_1e_1+\dots+y_ne_n)=\sum{a_{ij}x_iy_j}=x^tAy$,
где $x=(x_1,\dots,x_n)^T\in k^n$, $y=(y_1,\dots,y_n)^T\in F^n$~--- столбцы координат,
$a_{ij}=b(e_i,e_j)$, $A=(a_{ij})$~--- {\bf матрица Грама}.
Пусть $W$~--- подпространство $V$; определим {\bf ортогонал} к $W$:
$$
W^\perp=\{u\in V:b(u,w)=0\text{ для всех }w\in W\}.
$$
\end{definition}

\begin{lemma}\label{lemma:dimension_of_orthogonal}
$\dim W^\perp+\dim W\geq\dim V$.
\end{lemma}
\begin{proof}
Пусть $u_1,\dots,u_m$~--- базис $W$; построим линейное отображение $\a\colon V\to k^m$: $v\mapsto(b(v,u_i))_{i=1}^m$.
При этом $\Ker(\a)=W^\perp$, $\dim\Image(\a)\leq m=\dim W$, поэтому $\dim V=\dim\Ker(\a)+\dim\Image(\a)
\leq\dim W^\perp+\dim W$.
\end{proof}

Отображение
$\ph\colon V\to k$ называется {\bf квадратичным отображением} или
{\bf квадратичной формой},
и пара $(V,\ph)$ называется {\bf квадратичным пространством над $k$},
если $\ph$ удовлетворяет следующим условиям:
\begin{enumerate}
\item $\ph(av)=a^2\ph(v)$ для всех $a\in k$, $v\in V$;
\item отображение $b_\ph\colon V\times V\to k$, заданное формулой
$$
b_\ph(v,w):=\frac12(\ph(v+w)-\ph(v)-\ph(w)),
$$
является $k$-билинейным.
\end{enumerate}
При этом $b_\ph$ называется {\bf симметричной билинейной формой, ассоциированной с $\ph$}
(из определения очевидно, что $b_\ph$ симметрична).
Форма $\ph$ восстанавливается по $b_\ph$ формулой $\ph(v)=b_\ph(v,v)$.
Пусть $B=\{e_1,\dots,e_n)$~--- базис $V$. {\bf Матрицей [Грама] квадратичной формы} в базисе $B$
называется матрица $A=(b_\ph(e_i,e_j))_{\begin{smallmatrix}1\leq i\leq n\\1\leq j\leq n\end{smallmatrix}}$.
Легко видеть, что эта матрица симметрична. Обратно, по любой симметричной матрице из $M(n,k)$
строится квадратичная форма на $k^n$.
Если $x$~--- столбец координат некоторого вектора $v\in V$, то значение квадратичной формы на этом
векторе записывается так:
$$
\ph(v)=x^tAx.
$$
Значение билинейной симметричной формы $b_\ph$ на двух векторах $v,w\in V$ с координатными столбцами
$x$ и $y$ соответственно записывается так:
$$
b_\ph(v,w)=x^tAy=y^tAx.
$$

Два $n$-мерных векторых квадратичных пространства $(V,\ph)$ и $(V',\ph')$ называются
{\bf изометричными}, если существует $k$-линейный изоморфизм $T\colon V\to V'$
такой, что
$$
\ph(v)=\ph'(Tv)\text{ для всех $v\in V$.}
$$
Обозначение: $(V,\ph)\cong(V',\ph')$. В большинстве случаев мы забываем про пространства,
на которых определены формы, и пишем $\ph\cong\ph'$.
Очевидно, что изометричность является отношением эквивалентности.
Если в каждом из пространств $V,V'$ выбраны базисы, их можно отождествить с $k^n$,
и изоморфизм $T$ превращается в автоморфизм $k^n$, то есть, записывается
матрицей из $\GL(n,k)$. При этом если $A$~--- матрица $\ph$, $A'$~--- матрица $\ph'$,
то $x^tAy=(Tx)^tA'(Ty)$ для всех $x,y\in k^n$, откуда $A=T^tA'T$.

{\bf Определителем} $\ph$ называется определитель матрицы Грама $\ph$. Заметим,
что при замене базиса определитель матрицы Грама умножается на квадрат определителя
матрицы замены базиса; поэтому $\det(\ph)\in k^*/(k^*)^2\cup\{0\}$~--- определен
только с точностью до домножения на квадраты в поле $k$.

Пусть $(V_1,\ph_1)$, $(V_2,\ph_2)$~--- два квадратичных пространства над $k$ размерностей
$n_1$ и $n_2$ соответственно. По ним можно построить квадратичное пространство
$(V,\ph)$ размерности $n=n_1+n_2$:
$$
\begin{aligned}
V&=V_1\oplus V_2,\\
\ph(v)&=\ph_1(v_1)+\ph_2(v_2)
\end{aligned}
$$
для $v_1\in V_1$, $v_2\in V_2$, $v=v_1+v_2\in V$. Это пространство $(V,\ph)$ называется
{\bf прямой суммой} $(V_1,\ph_1)$ и $(V_2,\ph_2)$. Обозначается это так:
$(V,\ph)=(V_1,\ph_1)\oplus(V_2,\ph_2)$ или $(V_1,\ph_1)\perp(v_2,\ph_2)$.
Мы будем также писать
$\ph=\ph_1\oplus\ph_2=\ph_1\perp\ph_2$. Если $A_1$~--- матрица $\ph_1$, $A_2$~--- матрица
$\ph_2$ в некоторых базисах $B_1,B_2$ пространств $V_1,V_2$, то $B=B_1\sqcup B_2$~---
базис $V$, в котором $\ph$ имеет матрицу
$$
A=\begin{pmatrix} A_1 & 0 \\ 0 & A_2\end{pmatrix}.
$$
Аналогично можно определить сумму любого натурального количества квадратичных пространств.
Класс изометричности суммы зависит только от слагаемых, но не от их порядка.
Обратно, пусть $(V,\ph)$~--- квадратичное пространство и $\{V_i\}_{1\leq i\leq r}$~---
набор подпространств $V$ таких, что $V=V_1\oplus\dots\oplus V_r$ и $b_\ph(v_i,v_j)=0$
для всех $v_i\in V_i$, $v_j\in V_j$, $i\neq j$. Тогда $\ph=\ph_1\oplus\dots\oplus\ph_r$
для $\ph_i=\ph|_{V_i}$.

\begin{theorem}\label{thm:diagonalisation}
Любое квадратичное пространство $(V,\ph)$ над $k$ изометрично прямой сумме одномерных подпространств.
Другими словами, каждая $n$-арная квадратичная форма $\ph$ над $k$ эквивалентна диагональной форме
$\psi$ вида $\psi(x)=\sum_{i=1}^n{a_ix_i^2}$, $a_i\in k$.
\end{theorem}
\begin{proof}
Индукция по $n=\dim V$. Если $\ph(v)=0$ для всех $v\in V$, то $b_\ph=0$, и любой базис $V$ является
ортогональным.
Если $\ph(v_1)=a_1\neq 0$ для некоторого $v_1\in V$, рассмотрим подпространство
$$
U=(kv_1)^\perp=\{u\in V:b_\ph(u,v_1)=0\}
$$
всех векторов, ортогональных к $v_1$ (относительно $b_\ph$). При этом
по лемме~\ref{lemma:dimension_of_orthogonal} размерность $U$ не меньше,
чем $n-1$, но $v_1\notin U$,
поэтому $\dim U=n-1$, откуда $V=kv_1\oplus U$ и $\ph=\ph_1\oplus\ph_2$ для
$\ph_1=\ph|_{kv_1}$, $\ph_2=\ph|_{kv_2}$.
\end{proof}
Заметим, что в качестве $a_1$ можно взять любой элемент из $k^*$ вида $\ph(v_1)$ для $v_1\in V$.
\begin{proof}[{Второе доказательство.}] Приведем явный алгоритм. Будем действовать
индукцией по $n$; база $n=1$ очевидна. Пусть теперь $n>1$. Запишем нашу
форму в координатах с помощью какого-нибудь базиса $V$:
$\ph(x_1,\dots,x_n)=\sum_{i,j=1}^na_{ij}x_ix_j$. Предположим сначала,
что найдется диагональный коэффициент $a_{ii}\neq 0$. После перестановки базисных
векторов можно считать, что $a_{11}\neq 0$.
Посмотрим на слагаемые, содержащие $x_1$:
$\ph(x_1,\dots,x_n)=a_{11}x_1^2+2a_{12}x_1x_2+\dots+2a_{1n}x_1x_n+\ph'(x_2,\dots,x_n)$.
Выделим полный квадрат: $\ph(x_1,\dots,x_n)=a_{11}(x_1+\frac{a_{12}}{a_{11}}x_2+\dots
+\frac{a_{1n}}{a_{11}}a_n)^2+\ph''(x_2,\dots,x_n)$, и по предположению индукции
форма $\ph''$ от меньшего количества переменных приводится к диагональному виду.

Теперь предположим, что все диагональные коэффициенты равны 0, но найдется
недиагональный коэффициент $a_{ij}\neq 0$, $i\neq j$. После перестановки базисных
векторов можно считать, что $a_{12}\neq 0$ (а все $a_{ii}$ равны 0).
Сделаем замену: $x'_1=x_1+x_2$, $x'_2=x_1-x_2$. При этом
$\ph(x_1,\dots,x_n)=2a_{12}x_1x_2+\ph'(x_1,\dots,x_n)=\frac{1}{2}a_{12}{x'_1}^2-
\frac{1}{2}a_{12}{x'_2}^2+\ph''(x'_1,x'_2,x_3,\dots,x_n)$.
При этом $\ph''(x'_1,x'_2,x_3,\dots,x_n)$ не содержит мономов вида ${x'_1}^2$,
поскольку $\ph'(x_1,\dots,x_n)$ не содержит мономов вида $x_1x_2$, $x_1^2$ и $x_2^2$.
Значит, в новом базисе у нашей формы появился ненулевой диагональный коэффициент,
и можно выделить полный квадрат, как и выше.

Наконец, если все коэффициенты $\ph$ равны нулю, то форма нулевая и она уже
записана в диагональном виде.
\end{proof}
Диагональную форму $\ph(x)=\sum_{i=1}^n{a_ix_i^2}$ мы будем обозначать
$$
\ph=\la a_1,\dots,a_n\ra=\la a_1\ra\perp\dots\perp\la a_n\ra.
$$

Пусть $(V,\ph)$~--- квадратичное пространство, $A$~--- матрица формы $\ph$. Подпространство
$\rad V=V^\perp=\{u\in V:b_\ph(u,v)=0\text{ для всех } v\in V\}$ называется {\bf радикалом}
$(V,\ph)$. Пространство $(V,\ph)$ называется {\bf регулярным} или {\bf невырожденным},
если $\rad V=0$. Как всегда, мы часто говорим о регулярности (невырожденности) {\it формы},
опуская упоминание о пространстве.

Нетрудно видеть, что $\rad V=\{u\in V:u^tAv=0\text{ для всех } v\in V\}=\{u\in V:u^tA=0\}$;
поэтому $\rad V=0\Longleftrightarrow \det A\neq 0$; радикал и регулярность инвариантны относительно
изометрии; если $\ph$ не регулярно, то $\ph\cong\la a_1,\dots,a_{n-1},0\ra$, то есть
$\ph$ эквивалентна форме, зависящей лишь от $n-1$ переменных. Поэтому можно предполагать,
что все формы регулярны.
Более точно:
\begin{theorem}[о выделении регулярной части]\label{thm:regular_part}
Пусть $(V,\ph)$~--- квадратичная форма. Существует разложение $(V,\ph)=(W_0,\ph_0)\perp(W_1,\ph_1)$,
где $\ph_0(W_0)=0$ для всех $w_0\in W_0$, а $(W_1,\ph_1)$~--- невырожденная форма. Более того,
это разложение единственно с точностью до изометрии.
\end{theorem}
\begin{proof}
Существование такого разложения следует из теоремы~\ref{thm:diagonalisation}.
Заметим, что в любом подобном рзаложении $W_0\perp W_0$ и $W_0\perp W_1$, откуда
$W_0\perp V$, то есть $W_0\subset\rad(V)$. Если при этом $W_0\neq\rad(V)$, то
$\rad(V)\cap W_1\neq 0$, то есть в $W_1$ найдется вектор, ортогональный $V$,
чего не может быть по невырожденности $W_1$. Значит, $W_0=\rad(V)$.
Возьмем теперь два таких разложения:
$V=W_0\oplus W_1=W'_0\oplus W'_1$, при этом $W_0=W'_0=\rad(V)$.
Определим отображение $T\colon W_1\to W'_1$ как композицию
вложения $W_1\subset V$ и проекции $V$ на $W'_1$. По построению $T$
линейно, при этом для $w\in W_1$ разность $Tw-w$ лежит в $W_0=\rad(V)$.
Поэтому $\ph(Tw)=\ph(w+(Tw-w))=\ph(w)+2b_\ph(w,Tw-w)+\ph(Tw-w)$ и
два последних слагаемых равны 0. Значит, $T$~--- изометрия.
Заметим также, что $T$ можно продолжить до изометрии всего пространства,
если дополнить ее тождественным отображением на $W_0$.
\end{proof}

Пусть $\ph$~--- квадратичная форма над $k$, $L\supset k$~--- расширение полей. Тогда
$\ph$ можно рассматривать как квадратичную форму над $L$, которую мы будем обозначать
$\ph_L$ или $\ph\otimes L$. При этом
$$
\ph=\ph_K\text{ регулярна}\Longleftrightarrow \ph_L\text{ регулярна}.
$$

Пусть $(V,\ph)$~--- $n$-мерное квадратичное пространство над $k$
\begin{enumerate}
\item Для $a\in k$ будем говорить, что $\ph$ {\bf представляет $a$ над $k$}, если
существует ненулевой вектор $v\in V$ такой, что $\ph(v)=a$.
\item $\tilde{D_k}(\ph)=\{\ph(v):0\neq v\in V\}$~--- множество элементов $k$, представимых формой $\ph$.
\item $D_k(\ph)=D_k(\ph)\setminus\{0\}\subseteq k^*$.
\item $\ph$ называется {\bf универсальной (над $k$)}, если $D_k(\ph)=k^*$.
\item $\ph$ называется {\bf изотропной (над $k$)}, если $0\in\tilde{D_k}(\ph)$, иначе $\ph$
называется {\bf анизотропной (над $k$)}.
\end{enumerate}

\begin{example}
$x_1^2+x_2^2$ не универсальна, анизотропна над $\mathbb R$, но универсальная, изотропна
над $\mathbb C$.
\end{example}

Очевидно, что одномерное регулярное пространство не может быть изотропным. Посмотрим на двумерные.
\begin{proposition}\label{prop:hyperbolic_plane}
Есть только одна (с точностью до изометрии) регулярная изотропная квадратичная форма $\ph$ размерности 2,
а именно, $\ph(x)=2x_1x_2$. Кроме того, $\ph\cong\la a,-a\ra$ для любого $a\in k^*$. В частности,
$\ph$ универсальна.
\end{proposition}
\begin{proof}
Пусть $\ph$~--- двумерная регулярная изотропная форма на пространстве $V$ и $v_1\in V$, $\ph(v_1)=0$.
Поскольку $\ph$ регулярна, найдется $w\in V$ такой, что $b_\ph(v_1,w)\neq 0$. Домножая $w$ на подходящий
элемент $k^*$, можно считать, что $b_\ph(v_1,w)=1$. Для любого $\lambda\in k$ векторы
$v_1$ и $v_2=w+\lambda v_1$ образуют базис пространства $V$, в котором $\ph(v_1)=0$
и $b_\ph(v_1,v_2)=b_\ph(v_1,w+\lambda v_1)=b_\ph(v_1,w)+\lambda b_\ph(v_1,v_1)=1$.
Наконец, $\ph(v_2)=\ph(w+\lambda v_1)=\ph(w)+2\lambda b_\ph(w,v_1)+\lambda^2\ph(v_1)=
\ph(w)+2\lambda$. Значит, если положить $\lambda=-\ph(w)/2$, получим $\ph(v_2)=0$.
\end{proof}
Класс изометричности этой формы обозначается $\mathbb H\cong\la 1,-1\ra$ и называется
{\bf гиперболической плоскостью}.
Заметим, что $\det(\la 1,-1\ra)=-1$. Обратно, если $(V,\ph)$~--- двумерное квадратичное пространство
и $\det(\ph)=-1$, то $(V,\ph)$~--- гиперболическая плоскость.
\begin{proposition}\label{prop:isotropic_contains_hyperbolic_plane}
Пусть $(V,\ph)$~--- регулярное изотропное квадратичное пространство над $k$, $\dim V=n\geq 2$. Тогда
$V=U\oplus W$ и $U\cong\mathbb H$, $\dim W=n-2$, $\ph\cong\la 1,-1\ra\oplus\psi$, где $\psi=\ph|_W$.
\end{proposition}
\begin{proof}
Как и в предыдущем предложении, можно найти $v_1,v_2\in V$ такие, что двумерное подпространство
$U=kv_1+kv_2\subseteq V$ вместе с квадратичной формой $\ph|_U$ изоморфно гиперболической плоскости $\mathbb xH$.
Положим $W=U^\perp$, тогда $\dim W\geq n-2$ и $U\cap U^\perp=\rad U=0$, поскольку $U$ регулярно.
Значит, $\dim W=n-2$ и $V=U\oplus W$.
\end{proof}

\begin{theorem}
Для невырожденной формы $\ph$ и $a\in k^*$ равносильны:
\begin{enumerate}
\item $a\in D_k(\ph)$;
\item $\ph\perp\la -a\ra$ изотропна;
\item $\ph=\la a\ra\perp\ph_1$.
\end{enumerate}
\end{theorem}
\begin{proof}
$(1)\Rightarrow(3)$ из замечания после доказательства теоремы~\ref{thm:diagonalisation},
$(3)\Rightarrow(2)$ из предложения~\ref{prop:hyperbolic_plane},
$(2)\Rightarrow(1)$: если $V$~--- пространство формы $\ph$, то изотропность
$\ph\perp\la -a\ra$ на пространстве $V\perp kv_1$ означает, что для
некоторых $v\in V$, $\lambda\in k$, не равных одновременно 0,
выполняется $\ph(v)-a\lambda^2=0$. Если $v=0$, то $\lambda=0$; значит, $v\neq 0$.
Поэтому $\ph(v/\lambda)=a\lambda^2/\lambda^2=a$, что и требовалось.
\end{proof}

\begin{lemma}\label{lemma:representation_of_c}
Если форма $\la a,b\ra$ представляет элемент $c\in k^*$, то $\la a,b\ra\cong\la c,abc\ra$.
\end{lemma}
\begin{proof}
Из замечания после доказательства теоремы~\ref{thm:diagonalisation} ясно, что
$\la a,b\ra\cong\la c,d\ra$ для некоторого $d\in k$. Из сравнения определителей видно,
что $ab=cd$, поэтому $abc=c^2d$ и заменой второго базисного вектора формы $\la c,d\ra$ на
пропорциональный можно заменить $d$ на $abc$.
\end{proof}

%%================================================================================
\subsection{Теорема Витта о сокращении}

Пусть $(V,\ph)$~--- квадратичная форма; $v$~--- анизотропный вектор. Определим отражение $s_v$
относительно вектора $v$ формулой
$$
s_v(u)=u-2\frac{\ph(u,v)}{\ph(v,v)}v.
$$
Простое вычисление показывает, что отражение является изометрией.
\begin{lemma}
Пусть $v_1,v_2\in V$ и $\ph(v_1)=\ph(v_2)\neq 0$. Тогда существует композиция отражений, 
переводящая $v_1$ в $v_2$.
\end{lemma}
\begin{proof}
Если $\ph(v_1-v_2)\neq 0$, то подойдет отражение относительно $v_1-v_2$: $s_{v_1-v_2}(v_1)=v_2$.
Если $\ph(v_1+v_2)\neq 0$, то подойдет композиция отражения относительно $v_1+v_2$ ($s_{v_1+v_2}(v_1)=-v_2$)
и отражения относительно $v_2$.
Если же $\ph(v_1-v_2)=\ph(v_1+v_2)=0$, то $\ph(v_1)=\frac14\ph((v_1+v_2)+(v_1-v_2))=\frac12\ph(v_1+v_2,v_1-v_2)$
и $\ph(v_2)=\frac14\ph((v_1+v_2)-(v_1-v_2))=-\frac12\ph(v_1+v_2,v_1-v_2)$, откуда $\ph(v_1)=\ph(v_2)=0$,
что невозможно.
\end{proof}

\begin{corollary}
Любая изометрия невырожденного пространства есть композиция отражений.
\end{corollary}
\begin{proof}
Пусть $T:V\to V$ -- изометрия невырожденного квадратичного пространства $(V,\ph)$.
Доказываем индукцией по $n=\dim V$; база $n=1$ очевидна. Пусть $n>1$.
Возьмем $v\in V$ такой, что $\ph(Tv)=\ph(v)\neq 0$. По лемме найдется
композиция отражений $S\colon V\to V$ такая, что $Sv=Tv$. Отображение $S^{-1}T$, таким образом,
является изометрией и оставляет $v$ на месте; значит, $S^{-1}T$ оставляет на месте и $W=(kv)^\perp$~---
подпространство размерности $n-1$. По предположению индукции изометрия $S^{-1}T|_W$ является композицией
отражений (относительно векторов из $W$). Заметим, что любое отражение относительно вектора из $W$
оставляет на месте $v$, поскольку $v\perp W$. Значит, изометрия $S^{-1}T$ является композицией
тех же самых отражений, рассматриваемых уже как преобразований всего пространства $V$.
Перенося $S$ в другую часть, получаем, что и $T$ является композицией отражений.
\end{proof}

\begin{theorem}[Витта о сокращении]
Если $q\perp\ph_1\cong q\perp\ph_2$, то $\ph_1\cong\ph_2$.
\end{theorem}
\begin{proof}
Можно считать, что формы невырожденны; $q=\la a_1,\dots,a_n\ra$. Докажем, что из
$\la a\ra\perp\ph_1\cong\la a\ra\perp\ph_2$ следует, что $\ph_1\cong\ph_2$.
Пусть форма $\psi_1=\la a\ra\perp\ph_1$ задана на пространстве $kv_1\oplus W_1$,
а $\psi_2=\la a\ra\perp\ph_2$~--- на пространстве $kv_2\oplus W_2$. Изометричность
этих форм означает, что существует линейное отображение $T\colon kv_1\oplus W_1\to kv_2\oplus W_2$,
для которого $\psi_2(Tv)=\psi_1(v)$.
%При этом $v=\alpha v_1+w_1$, $w_1\in W_1$ и $\psi_1(v_1)=a$, $b_{psi_1}(v_1,w_1)=0$.
Запишем $Tv_1=xv_2+w_2$. Тогда
$\psi_2(v_2)=a$ и $\psi_2(Tv_1)=a$. По лемме найдется изометрия $S\colon kv_2\oplus W_2\to kv_2\oplus W_2$
такая, что $Sv_2=Tv_1$. Рассмотрим отображение $S^{-1}T\colon kv_1\oplus W_1\to kv_2\oplus W_2$.
Нетрудно видеть, что $S^{-1}T$ является изометрией между $\psi_1$ и $\psi_2$;
кроме того, $S^{-1}Tv_1=v_2$, поэтому $S^{-1}T$ переводит $W_1=(kv_1)^\perp$ в $W_2=(kv_2)^\perp$.
Это означает, что ограничение $S^{-1}T$ на $W_1$ и дает нужную изометрию между $\ph_1$ и $\ph_2$.
\end{proof}

\begin{corollary}[{\bf Теорема о продолжении изометрии}]\label{cor:isometry-extension}
Пусть $(V,\ph)$~--- квадратичное пространство, $W_1,W_2$~--- подпространства в $V$,
и пусть $\a\colon W_1\to W_2$~--- изометрия между ними такая,
что $\a(W_1\cap\rad V) = W_2\cap\rad V$.
Тогда существует изометрия $\b\colon V\to V$ такая, что $\b|_{W_1}=\a$.
\end{corollary}

\begin{proof}
Сначала докажем теорему для случая $\rad V=0$ (то есть, форма $\ph$ невырождена на $V$).

В случае, когда $W_1$ невырождено, утверждение следует из теоремы Витта о сокращении.
Действительно, в этом случае $V$ раскладывается в прямую сумму $W_i$ и его ортогонального
дополнения ($i=1,2$). По теореме о сокращении существует изометрия 
$\gamma:W_1^\perp\to W_2^\perp$. Тогда $\b=(\a,\gamma)$~--- изометрия $V\to V$.

Если же $W_1$~--- вырожденное подпространство в невырожденном пространстве
$V$, то можно выбрать базис $u_1,\dots,u_{2m},v_1,\dots,v_k$ пространства $V$, содержащий базис 
$u_1,u_3,\dots,u_{2s-1}$,\allowbreak $v_1,\dots,v_r$ пространства $W_1$ такой, 
что матрица формы $\ph$ в этом базисе будет иметь вид
$$
\operatorname{diag}\left(\left(\begin{smallmatrix}0&1\\1&0\end{smallmatrix}\right),\dots,
\left(\begin{smallmatrix}0&1\\1&0\end{smallmatrix}\right),\alpha_1,\dots,\alpha_k\right).
$$
В этом случае нетрудно распространить изометрию на невырожденное подпространство, порожденное
$u_1,\dots,u_{2s},v_1,\dots,v_r$, а затем использовать теорему о сокращении.

Перейдем к общему случаю.
Выберем подпространство $\tilde W_1$ такое, что
$W_1 = (W_1\cap\rad V)\oplus\tilde W_1$,
и расширим его до подпространства $\tilde V$ такого,
что $V = \rad V \oplus\tilde V$.
Тогда
$$
W_2 = \a(W_1) = \a(W_1\cap\rad V)\oplus \a(\tilde W_1) = (W_2\cap\rad V)\oplus \a(\tilde W_1).
$$
Расширим $\a(\tilde W_1)$ до подпространства $\tilde V'$
так, что $V = \rad V\oplus\tilde V'$.
Заметим, что теперь у нас есть два ортогональных дополнения
до $\rad V$: $\tilde V$ и $\tilde V'$.
По теореме о выделении регулярной части, невырожденные
подпространства $\tilde V$ и $\tilde V'$ изометричны.
К их подпространствам $\tilde W_1\leq\tilde V$
и $\a(\tilde W_1)\leq\tilde V'$ можно применить только что
доказанный случай невырожденный формы
и получить изометрию $\tilde V\to \tilde V'$,
продолжающую $\a_{\tilde W_1}$.
Далее, изометрию $\a|_{W_1\cap\rad V}$ между $W_1\cap\rad V$
и $\a(W_1\cap\rad V)$ можно продолжить до изометрии
$\rad V\to\rad V$ произвольным линейным отображением,
поскольку форма на этих подпространствах нулевая.
Ортогональная прямая сумма этих двух изометрий
дает нужную изометрию на $V$.
\end{proof}

\begin{corollary}
Любая невырожденная форма $\ph$ представляется в виде
$$\ph\cong\underbrace{\mathbb H\perp\dots\perp\mathbb H}_{r\text{ раз}}\perp\ph_{an},$$
где {\bf анизотропная часть} $\ph_{an}$
определена однозначно с точностью до изометрии, и {\bf индекс Витта} $i(\ph):=r$ определен однозначно.
\end{corollary}
\begin{proof}
По предложению~\ref{prop:isotropic_contains_hyperbolic_plane} если форма изотропна, из нее можно выделить
$\mathbb H$. Продолжая этот процесс, дойдем до какой-то анизотропной формы (потому что размерность все
время убывает). Осталось проверить единственность. Предположим, что
$\ph\cong\perp\ph_{an}$,
$\ph\cong\underbrace{\mathbb H\perp\dots\perp\mathbb H}_{r\text{ раз}}\perp \psi
\cong\underbrace{\mathbb H\perp\dots\perp\mathbb H}_{r'\text{ раз}}\perp\psi'$, где $\psi,\psi'$ анизотропны.
Не умаляя общности, $r\geq r'$. Если $r>r'$, то сокращая (по теореме Витта) слева и справа $r'$
раз на $\mathbb H$,
получаем, что $\underbrace{\mathbb H\perp\dots\perp\mathbb H}_{r-r'\text{ раз}}\perp\psi\cong\psi'$,
но слева стоит изотропная форма, а справа~--- анизотропная.
Поэтому $r=r'$ и после сокращения получаем $\psi\cong\psi'$.
\end{proof}

Пока что считаем все квадратичные формы невырожденными и диагональными.
Обозначим $G(k)=k^*/(k^*)^2$~--- {\bf square class group}.
Пусть $\ph\cong\la a_1,\dots,a_m\ra$, $a_i\in k^*$;
$\det\ph=(\prod a_i)(k^*)^2\in G(k)$~--- {\bf определитель (детерминант)} $\ph$,
$d(\ph)=(-1)^{m(m-1)/2}\det\ph\in G(k)$~--- {\bf дискриминант} $\ph$.
Если $\ph\cong\la a_1,\dots,a_m\ra$, $\psi\cong\la b_1,\dots,b_n\ra$~--- две формы, то
$\ph\perp\psi\cong\la a_1,\dots,a_m,b_1,\dots,b_n\ra$~--- {\bf (ортогональная) сумма} $\ph$ и $\psi$,
$\ph\otimes\psi\cong\la\dots,a_ib_j,\dots\ra_{\begin{smallmatrix}1\leq i\leq m\\1\leq j\leq n\end{smallmatrix}}$~---
{\bf (тензорное) произведение} $\ph$ и $\psi$.
Если $\ph\cong\la a_1,\dots,a_m\ra$ и $a\in k^*$, то $a\ph\cong\la aa_1,\dots,aa_m\ra$~--- произведение
$\ph$ на $a$. Если $r\in\mathbb N$, то $r\times\ph\cong\underbrace{\ph\perp\dots\perp\ph}_{r\text{ раз}}$~---
$r$-кратная сумма $\ph$ с собой, $0\times\ph=0$~--- пустая форма размерности 0.

Рассмотрим абелев моноид невырожденных квадратичных форм относительно ортогональной суммы;
по теореме Витта он является моноид с сокращением, поэтому он вкладывается в свою группу
%??? подробнее?
Гротендика. На этой абелевой группе определено умножение, индуцированной тензорным
произведением. В результате получаем {\bf кольцо} (коммутативное, ассоциативное, с 1)
{\bf Витта--Гротендика} $\tilde W(k)$.

\begin{definition}
Квадратичная форма $\ph$ называется {\bf гиперболической}, если она изоморфна
прямой сумме гиперболических плоскостей: $\ph=r\times\mathbb H$ для некоторого $r\geq 0$.
\end{definition}
\begin{proposition}
Гиперболические формы (и противоположные к ним) образуют идеал в кольце $\tilde W(k)$.
\end{proposition}
\begin{proof}
Очевидно, что сумма гиперболических форм гиперболична; поскольку одномерные формы аддитивно
порождают $\tilde W(k)$, достаточно доказать, что $\mathbb H\otimes\la a\ra$ гиперболична,
но $\mathbb H\otimes\la a\ra\cong\la a,-a\ra\cong\mathbb H$.
\end{proof}
\begin{definition}
Фактор-кольцо кольца Витта--Гротендика $\tilde W(k)$ по идеалу гиперболических форм
называется {\bf кольцом Витта} и обозначается $W(k)$.
\end{definition}

Вот другое определение кольца Витта: две невырожденные
формы $\ph,\psi$ над $k$ называются {\bf подобными}, если $\ph_{an}\cong\psi_{an}$.
Обозначение: $\ph\sim\psi$. Множество классов эквивалентности регулярных форм над $k$
обозначается $W(k)$.
\begin{theorem}[Витт, 1937]
Множество $W(k)$ является коммутативным ассоциативным кольцом с 1
относительно операций, индуцированных $\oplus$ и $\otimes$ и называется {\bf кольцом Витта}.
Относительно операции $\oplus$ множество $W(k)$ является абелевой группой и называется {\bf группой Витта}.
\end{theorem}
\begin{proof}
Очевидно.
\end{proof}

Заметим, что в качестве представителя элемента $W(k)$ можно взять квадратичную форму (а не формальную разность
двух квадратичных форм, как в $\tilde W(k)$), и для ненулевого класса эту форму можно выбрать анизотропной.

\begin{examples}
\begin{enumerate}
\item Если поле $k$ алгебраически замкнуто, то $W(k)=\mathbb Z/2\mathbb Z$.
\item $W(\mathbb R)=\mathbb Z$.
\item $W(\mathbb Z/p\mathbb Z)=\mathbb Z/2\mathbb Z\oplus\mathbb Z/2\mathbb Z$ для $p\equiv 1\pmod 4$ и
$W(\mathbb Z/p\mathbb Z)=\mathbb Z/4\mathbb Z$ для $p\equiv 3\pmod 4$.
\end{enumerate}
\end{examples}

% 8.03.2010

Дадим еще одну характеризацию индекса Витта.

\begin{definition}
Пусть $(V,q)$~--- квадратичная форма. Подпространство $W\leq V$ называется {\bf вполне изотропным},
если $q|_W=0$. Это условие равносильно тому, что всякий вектор $v\in W$ изотропен.
\end{definition}

\begin{proposition}\label{prop:totally_isotropic_subspaces}
Пусть $(V,q)$~--- квадратичная форма. Все максимальные вполне изотропные подпространства $V$
имеют одинаковую размерность, равную индексу Витта $i(q)$ формы $q$.
\end{proposition}
\begin{proof}
Очевидно, что если $q=(m\times\mathbb H)\perp\ph$, то в $V$ есть вполне изотропное подпространство
размерности $m$. Обратно, пусть в $V$ есть такое подпространство. Докажем индукцией по $m$, что
тогда в $V$ вкладывается сумма $m$ гиперболических плоскостей. При $m=0$ доказывать нечего.
Если $m>0$, выберем изотропный вектор $v\in V$. Рассуждение из
доказательства предложения~\ref{prop:isotropic_contains_hyperbolic_plane} показывает, что 
найдется вектор $v'\in V$ такой, что $kv\oplus kv'\cong\mathbb H$; стало быть, $q\cong\mathbb H\perp q'$.
Пусть $W=v^\perp$. Тогда $kv\subset W$ и на факторе $W/kv$ возникает корректно определенная
форма $\tilde q$, задаваемая равенством $\tilde q(w+kv):=q(w)$. Легко видеть, что $\tilde q\cong q'$.
Но по построению $\tilde q$ имеет вполне изотропное подпространство размерности $m-1$,
поэтому такое подпространство есть и в $q'$. По предположению индукции, в $q'$ есть сумма $m-1$
гиперболических плоскостей, поэтому в $q$ есть сумма $m$ гиперболических плоскостей.
\end{proof}

\begin{lemma}\label{lem:isotropicdiff}
Пусть $q, q'$~--- две анизотропные формы. Предположим, что $i(q\perp -q')\geq n$. Тогда существуют
квадратичные формы $\ph, q_1, q'_1$ такие, что $\dim\ph=n$ и $q\cong\ph\perp q_1$ и $q'\cong\ph\perp q'_1$.
\end{lemma}
\begin{proof}
Индукция по $n$. Пусть $n=1$: $q\perp -q'$ изотропна, поэтому найдутся $x\in V$, $x'\in V'$ такие, что
$q(x)=q'(x')\neq 0$ (здесь $V, V'$~--- подлежащие пространства форм $q$ и $q'$ соответственно), 
и утверждение очевидно. Если $n>1$, действуя так же, получаем, что $q\cong\la a\ra\perp q_2$,
$q'\cong\la a\ra\perp q'_2$ для некоторых $a,q_2,q'_2$. Тогда $i(q_2\perp q'_2)\geq n-1$ и
можно применить индукционное предположение.
\end{proof}

\begin{lemma}\label{lem:abaplusb}
Пусть $a,b\in k^*$. Тогда
$\la a,b\ra\cong\la a+b,ab(a+b)\ra$.
\end{lemma}
\begin{proof}
Немедленно следует из леммы~\ref{lemma:representation_of_c}.
\end{proof}

\begin{theorem}\label{theorem:witt_ring_generators_and_relations}
\begin{enumerate}
\item Аддитивная группа кольца $\tilde W(k)$ порождается (как абелева группа) образующими $\la a\ra$, $a\in k^*$,
удовлетворяющими соотношениям $\la ab^2\ra=\la a\ra$ и $\la a,b\ra=\la a+b,ab(a+b)\ra$.
\item Аддитивная группа кольца $W(k)$ порождается (как абелева группа) образующими $\la a\ra$, $a\in k^*$,
удовлетворяющими соотношениям $\la ab^2\ra=\la a\ra$, $\la a,b\ra=\la a+b,ab(a+b)\ra$ 
и дополнительным соотношениям $\la -a\ra=-\la a\ra$.
\end{enumerate}
\end{theorem}
\begin{proof}
Пусть $V(k)$~--- группа, порожденная соотношениями из первого пункта формулировки теоремы. Обозначим
через $[a]$ образующую, соответствующую скаляру $a\in k^*$. Предыдущие результаты показывают, что
существует сюръективный гомоморфизм $V(k)\to\tilde W(k)$, переводящий $[a]$ в $\la a\ra$.
Для доказательства первого пункта остается показать, что если $a_1,\dots,a_n,b_1,\dots,b_n\in k^*$
таковы, что $\la a_1,\dots,a_n\ra\cong\la b_1,\dots,b_n\ra$, то
$[a_1]+\dots+[a_n]=[b_1]+\dots+[b_n]$.

Будем действовать индукцией по $n$ с тривиальной базой $n=1$. Пусть $n=2$. Поскольку
$\la a_1,a_2\ra\cong\la b_1,b_2\ra$, то найдутся $x_1,x_2\in k$ такие, что $b_1=a_1x_1^2+a_2x_2^2$.
Если $x_2=0$, то $\la b_1\ra=\la a_1\ra$, откуда $\la b_2\ra=\la a_2\ra$ и доказывать нечего.
Если $x_1=0$, все аналогично. Если же $x_1x_2\neq 0$, заменяя $a_i$ на $a_ix_i^2$, можно
считать, что $x_1=x_2=1$. По лемме~\ref{lem:abaplusb} имеем $\la b_2\ra\cong\la a_1a_2(a_1+a_2)\ra$
и доказательство окончено.

Наконец, предположим, что $n\geq 3$. Обозначим $q=\la a_1,\dots,a_{n-1}\ra$, $q'=\la b_1,\dots,b_{n-1}\ra$.
Тогда $q\perp -q'\sim\la b_n,-a_n\ra$, откуда, по лемме~\ref{lem:isotropicdiff} существуют
$c_1,\dots,c_{n-2},e,f\in k^*$ такие, что $q\cong\la c_1,\dots,c_{n-2},e\ra$ и
$q'\cong\la c_1,\dots,c_{n-1},f\ra$ и по теореме Витта $\la e,a_n\ra\cong\la f,b_n\ra$.
Применяя индукционное предположение, получаем $[a_1]+\dots+[a_{n-1}]=[c_1]+\dots+[c_{n-2}]+[e]$,
$[b_1]+\dots+[b_{n-1}]=[c_1]+\dots+[c_{n-2}]+[f]$ и $[e]+[a_n]=[f]+[b_n]$,
и отсюда все следует. Второй пункт теоремы доказывается совершенно аналогично.
\end{proof}

\subsection{Первая теорема Касселса о представимости}

Пусть $\ph$~--- квадратичная форма над $k$, $k(t)$~--- {\bf поле рациональных функций}
над $k$ от одной переменной $t$.

\begin{lemma}
Если $\ph$ анизотропна над $k$, то $\ph_{k(t)}$ анизотропна над $k(t)$.
\end{lemma}
\begin{proof}
Пусть $\ph(f)=0$, где $f=(f_1,\dots,f_n)$, $f_i\in k(t)$. Пусть $g_0$~---общий знаменатель
функций $f_i$: $f_i=g_i/g_0$, где $g_0,g_1,\dots,g_n\in k[t]$. Тогда $\ph(g)=g_0^2(f)=0$
для $0\neq g=(g_1,\dots,g_n)$. Теперь пусть $d=\gcd(g_1,\dots,g_n)\in k[t]$;
$g_i=dh_i$, $h_i\in k[t]$~--- взаимно просты. Пусть $h=(h_1,\dots,h_n)$, тогда
$\ph(g)=d^2\ph(h)=0$~--- тождество. Поскольку $k[t]$~--- область целостности, имеем
$\ph(h)=0$. Положим $c_i=h_i(0)\in k$, $c=(c_1,\dots,c_n)$~--- ненулевой вектор
(иначе все $h_i$ делились бы на $t$). Поэтому $c\in k^n$ и $\ph(c)=0$, противоречие.
\end{proof}
\begin{theorem}
Пусть $\ph(x)=\ph(x_1,\dots,x_n)=\sum_{i,j=1}^n{a_{ij}x_ix_j}$~--- $n$-арная квадратичная форма
над $k$. Пусть $0\neq p(t)\in k[t]$. Предположим, что $\ph$ представляет $p$ над полем $k(t)$.
Тогда $\ph$ представляет $p$ над кольцом $k[t]$, то есть найдутся $f_i\in k[t]$ такие,
что $\ph(f_1,\dots,f_n)=p$.
\end{theorem}
\begin{proof}
Если $\ph$ не регулярна, можно заменить $\ph$ на $(n-1)$-форму и действовать по индукции.
Если $n=1$, $\ph(x)=a_{11}x_1^2$, $a_{11}f_1^2=p$ для $f_1\in k(t)$, откуда $f_1\in k[t]$.
Предположим теперь, что $\ph$ регулярна, но изотропна. Тогда $\ph\cong H\perp\psi$ над $k$,
то есть можно считать, что $\ph(x)=2x_1x_2+\psi(x_3,\dots,x_n)$. Положим $x_1=p(t)$, $x_2=1/2$,
$x_3=\dots=x_n=0$.
Наконец, $\ph$ регулярна и анизотропна. По предположению
$$
\ph\left(\frac{f_1}{f_0},\dots,\frac{f_n}{f_0}\right)=p
$$
для некоторых многочленов $f_0,\dots,f_n\in k[t]$; не умаляя общности, имеем $\gcd(f_0,\dots,f_n)=1$.
Более того, можно считать, что из всех представлений в таком виде выбрано то, у которого
$d=\deg f_0$ минимальна. Предположим, что $d>0$ и получим противоречие.
Рассмотрим новую форму $\psi=\la -p(t)\ra\oplus\ph_{k(t)}$ над $k(t)$:
$\psi(x_0,\dots,x_n)=-p(t)x_0^2+\ph(x_1,\dots,x_n)$. Очевидно,
что $\psi(f_0,\dots,f_n)=0$.
Поделим с остатком $f_i$ на $f_0$: $f_i=f_0g_i+r_i$, $\deg r_i<d$.
В частности, $g_0=1$, $r_0=0$, $\deg r_0=-\infty$.
$\psi(g)\neq 0$ по минимальности $d=\deg f_0$. В частности,
$f$ и $g$ линейно независимы над $k(t)$.
Определим $h=\lambda f-\mu g\in (k(t))^{n+1}$, $\lambda=\psi(g)$,
$\mu=2b_\psi(f,g)$. $h=(h_0,\dots,h_n)$, $\lambda\neq 0$, значит, $h\neq 0$.
Но
$$
\psi(h)=\lambda^2\psi(f)-2\lambda\mu b_\psi(f,g)+\mu^2\psi(g)=0.
$$
На самом деле $h_0\neq 0$, иначе $\psi$ была бы изотропна над $k(t)$ и,
по лемме, над $k$. Осталось оценить $\deg h_0$:
$$
\begin{aligned}
h_0=\lambda f_0-\mu=\psi(g)f_0-b_\psi(f,g)&=\frac{1}{f_0}\psi(f_0g-f)\\
&=\frac{1}{f_0}\sum_{i,j=1}^na_{ij}(f_0g_i-f_i)(f_0g_j-f_j).\\
\end{aligned}
$$
Поэтому $\deg\psi(f_0g-f)\leq 2\max_{i=1,\dots,n}\deg(f_0g_i-f_i)=2\max_{i=1,\dots,n}\deg r_i
\leq 2(d-1)$,
откуда $\deg h_0=-d+\deg\psi(f_0g-f)\leq d-2$, противоречие.
\end{proof}

\begin{theorem}[обобщение]
Пусть $\ph(x)=\sum_{i,j=1}^na_{ij}x_ix_j$~--- квадратичная форма над $k(t)$ такая, что
$a_{ij}\in k[t]$ и $\deg a_{ij}\leq 1$ для всех $(i,j)$. Предположим, что $\ph$
анизотропна над $k(t)$. Пусть $\ph$ представляет над $k(t)$ многочлен $0\neq p(t)\in k[t]$.
Тогда $\ph$ представляет $p(t)$ над $k[t]$.
\end{theorem}
\begin{proof}
Доказательство повторяется, но в этот раз $\deg\psi(f_0g-f)\leq 1+2\max\deg r_i\leq 2d-1$,
откуда $\deg h_0\leq d-1<d$.
\end{proof}
\begin{remark}
Доказательство перестает быть верным, если $\ph$ изотропна! Пусть $\ph=\la t,-t\ra$, $p(t)=1$;
тогда $\ph$ представляет $p$ над $k(t)$, но не над $k[t]$.
\end{remark}

% 15.03.2010

\subsection{Теорема о подформе}

\begin{theorem}[Принцип подстановки]
Пусть $\ph$~--- $n$-арная квадратичная форма над $k$, $0\neq p=p(t_1,\dots,t_m)\in k[t_1,\dots,t_m]$
и $c_1,\dots,c_m\in k$. Если $\ph$ представляет $p$ над полем рациональных функций $k(t_1,\dots,t_m)$,
то $\ph$ представляет элемент $p(c_1,\dots,c_m)$ над $k$.
\end{theorem}
\begin{proof}
Индукция по $m$.
\end{proof}

\begin{lemma}
Пусть $d,a_1,\dots,a_n\in k^*$; предположим, что $\ph=\la a_1,\dots,a_n\ra$ представляет многочлен
$d+a_1t^2$ над $k(t)$. Тогда или $\ph$ изотропна над $k$ или $\ph'=\la a_2,\dots,a_n\ra$
представляет $d$ над $k$.
\end{lemma}
\begin{proof}
Предположим, что $\ph$ анизотропна.
Из теоремы Касселса получаем, что $\sum_{i=1}^n{a_if_i^2}=d+a_1t^2$ для некоторых $f_i\in k[t]$.
Легко видеть, что $\deg f_i=1$, пусть $f_i=b_i+c_it$; уравнение $b_1+c_1t=\pm t$ имеет некоторое
решение $t=c$. Подставляя $c$, видим, что $\ph'$ представляет~$d$.
\end{proof}

\begin{theorem}[Теорема о подформе]\label{theorem:subform}
Пусть $\ph\cong\la a_1,\dots,a_n\ra$, $\psi\cong\la b_1,\dots,b_m\ra$~--- регулярные квадратичные формы
над $k$. Предположим, что $\ph$ анизотропна. Следующие утверждения эквивалентны:
\begin{enumerate}
\item $\psi$ изоморфна подформе $\ph$, то есть $\ph\cong\psi\perp\xi$ для некоторой квадратичной формы
$\xi$ над $k$ (возможно, $\xi=0$).
\item $D_L(\psi)\subseteq D_L(\ph)$ для любого поля $L\supseteq k$.
\item $\ph$ представляет <<общее значение>> $\psi$, то есть $\ph$ представляет
$\psi(t_1,\dots,t_m)=b_1t_1^2+\dots+b_mt_m^2$ над полем рациональных функций $k(t_1,\dots,t_m)$.
\end{enumerate}
В частности, из любого из этих утверждений следует, что $m\leq n$.
\end{theorem}
\begin{proof}
$(1)\Rightarrow(2)\Rightarrow(3)$~--- очевидно. Докажем $(3)\Rightarrow(1)$ индукцией по $m$,
база тривиальна. Пусть теперь $m>0$. По принципу подстановки $\ph$ представляет $b_1\neq 0$ над $k$.
Значит, мы можем записать $\ph\cong\la b_1\ra\perp\ph'$, где $\ph'$ автоматически анизотропна.
Поскольку $\ph$ представляет $b_1t_1^2+(b_2t_2^2+\dots+b_mt_m^2)$ над $k(t_2,\dots,t_m)(t_1)$,
по лемме $\ph'$ представляет $d=\psi'(t_2,\dots,t_m)=b_2t_2^2+\dots+b_mt_m^2$.
Теперь можно применить предположение индукции к паре $(\ph',\psi')$ и получить,
что $\ph'\cong\psi'\perp\xi$ и $\ph\cong\la b_1\ra\perp\ph'\cong\la b_1\ra\perp\psi'\perp\xi\cong\psi\perp\xi$.
\end{proof}

\begin{definition}
В ситуации пункта 1 теоремы~\ref{theorem:subform} будем говорить, что $\psi$~--- {\bf подформа}
$\ph$ и писать $\psi\leq\ph$.
\end{definition}

\begin{definition}
Пусть $\ph$~--- квадратичная форма. Напомним, что $D(\ph)=\{a\in k^*\mid\exists x, \ph(x)=a\}$~--- множество
ненулевых элементов, представляемых формой $\ph$. Положим
$G(\ph)=\{a\in k^*\mid a\ph\cong\ph\}$~--- множество {\bf коэффициентов подобия} $\ph$.
\end{definition}
\begin{lemma}
\begin{enumerate}
\item Если $\ph\leq\ph'$, то $D(\ph)\subseteq D(\ph')$.
\item Если $\ph$ изотропна, то $D(\ph)=k^*$.
\item Для любого $\lambda\in k^*$ имеем $G(\lambda\ph)=G(\ph)$.
\item $G(\ph)$ зависит лишь от класса $\ph$ в кольце Витта $W(k)$.
\item $G(\ph)$~--- подгруппа $k^*$, содержащая $(k^*)^2$. Если $a\in G(\ph)$, $b\in D(\ph)$,
то $ab\in D(\ph)$.
\item Если $1\in D(\ph)$, то $G(\ph)\subseteq D(\ph)$.
\end{enumerate}
\end{lemma}
\begin{proof}
Пункты 1--3 очевидны, для доказательства 4 достаточно проверить, что $G(\ph)=G(\ph\perp\mathbb H)$.
Заметим, что $a\mathbb H\cong\la a,-a\ra\cong\mathbb H$ для любого $a\in k^*$.
Если $a\in G(\ph)$, то $\ph\cong a\ph$, поэтому
$$
\ph\perp\mathbb H\cong a\ph\perp\mathbb H\cong a\ph\perp a\mathbb H\cong a(\ph\perp\mathbb H).
$$
Обратно, если $\ph\perp\mathbb H\cong a(\ph\perp\mathbb H)\cong a\ph\perp a\mathbb H\cong
a\ph\perp\mathbb H$, то по теореме Витта о сокращении получаем, что $\ph\cong a\ph$.
Далее, 5 очевидно и 6 следует из 5.
\end{proof}

\begin{lemma}
Пусть $\ph$~--- квадратичная форма над $k$ и $\ph'\leq\ph$. Если $\dim\ph'>\dim\ph-i(\ph)$,
то $\ph'$ изотропна.
\end{lemma}
\begin{proof}[Первое доказательство]
Пусть $V$~--- пространство формы $\ph$, $W$~--- подпространство, соответствующее $\ph'$,
$H\leq V$~--- максимальное вполне изотропное подпространство размерности $i(\ph)$
(см.~предложение~\ref{prop:totally_isotropic_subspaces}).
При этом $\dim(W)+\dim(H)>\dim(V)$, откуда пересечение $W\cap H$ непусто.
\end{proof}
\begin{proof}[Второе доказательство]
Запишем $\ph'\perp\ph''\cong\ph\cong\ph_{an}\perp i(q)\times\mathbb H$ для некоторой формы $\ph''$.
Тогда $\ph'\perp\ph''\perp-\ph''\cong\ph_{an}\perp-\ph''\perp i(q)\times\mathbb H$. Заметим,
что $\ph''\perp-\ph''\cong\dim(\ph'')\times\mathbb H$, поэтому
$\ph'\cong\ph_{an}\perp-\ph''\perp(i(q)-\dim(\ph''))\times\mathbb H$ и $\ph'$ изотропна.
\end{proof}

\subsection{Поведение квадратичных форм при конечных расширениях полей}

Посмотрим на самый простой нетривиальный случай~--- квадратичное расширение.
\begin{lemma}
Пусть $L=k(\sqrt{a})$, $\ph$~--- анизотропная форма над $k$. Тогда равносильны:
\begin{enumerate}
\item $\ph_L$ изотропна;
\item $\ph=b\la 1,-a\ra\perp\psi$ для некоторых $b\in k^*$ и формы $\psi$.
\end{enumerate}
\end{lemma}
\begin{proof}
Очевидно, что из второго пункта следует первый. Пусть теперь $\ph_L$ изотропна. Это означает,
что в $V\otimes L$ есть изотропный вектор, то есть, $\ph(v+w\sqrt{a})=0$ для некоторых
$v,w\in V$, не равных одновременно нулю. Значит, $\ph(v)+a\ph(w)+2\sqrt{a}b_\ph(v,w)=0$,
откуда $\ph(v)=-a\ph(w)$ и $b_\ph(v,w)=0$. Разложим $V$ в прямую сумму
пространства $W=kv+kw$ и ортогонального дополнения $W^\perp$. Относительно этого разбиения
и получим необходимое разложение формы $\ph$.
\end{proof}

\begin{theorem}
Пусть $L=k(\sqrt{a})$, $\ph$~--- анизотропная форма над $k$. Тогда равносильны:
\begin{enumerate}
\item $i(\ph_L)\geq i$;
\item $\ph=\la 1,-a\ra\otimes\ph'\perp\psi$, где $\rk\ph'=i$.
\end{enumerate}
\end{theorem}
\begin{proof}
Индукция по $i$; используется, что $i(\ph\perp\mathbb H)=i(\ph)+1$.
\end{proof}

\begin{corollary}
В условиях теоремы если $\ph_L$ гиперболична, что $\ph=\la 1,-a\ra\otimes\psi$.
\end{corollary}

Таким образом, ядро отображения $W(k)\to W(k(\sqrt{a}))$ порождается формами вида $\la 1,-a\ra$.

\begin{theorem}[Спрингер]
Пусть $k$~--- подполе $L$ и степень $[L:k]$ нечетна. Тогда $i(\ph_L)=i(\ph)$.
\end{theorem}
\begin{proof}
Достаточно проверить, что из анизотропности $\ph$ следует анизотропность $\ph_L$ и проверить
при этом лишь случай расширения, порожденного одним элементом: $L=k(\alpha)$.
Будем доказывать индукцией по $n=[L:k]$ с тривиальной базой $n=1$.
Предположим противное: пусть $P$~--- минимальный многочлен $\alpha$, $d=\dim(\ph)$ и
$(x_1,\dots,x_d)\in L^d\setminus\{0\}$ таковы, что $\ph(x_1,\dots,x_d)=0$. Можно записать
$x_i=g_i(x_\alpha)$, где $g\in k[t]$, $m:=\max(\deg g_i)<n$ и $g_i$ не все равны 0.
Разделив на наибольший общий делитель, можно считать, что они взаимно просты в совокупности.
Получаем равенство в $k[t]$:
$$
\ph(g_1,\dots,g_d)=P\cdot h
$$
для некоторого $h\in k[t]$. При этом $\deg(h)=2m-n\leq n-2$: действительно, из анизотропности
$\ph$ следует, что $\deg(\ph(g_1,\dots,g_d))=2m$. В частности, $\deg(h)$ нечетна.
Пусть $h'$~--- неприводимый множитель $h$ нечетной степени; очевидно, что $\deg(h')\leq n-2$.
Обозначим $F=k[t]/(h')$; это расширение $k$ нечетной степени, меньшей $n$. Пусть $\beta$~---
образ $t$ в $F$. Тогда $\ph_F(g_1(\beta),\dots,g_d(\beta))=0$. По предположению индукции
$\ph_F$ анизотропна, поэтому $g_1(\beta)=\dots=g_d(\beta)=0$, откуда $h'$ является общим делителем
$g_1,\dots,g_d$, что противоречит предположению.
\end{proof}

\begin{corollary}
Если степень $L$ над $k$ нечетна, то отображение $W(k)\to W(L)$ инъективно.
\end{corollary}

\section{Теория Пфистера}


\subsection{Формы Пфистера}

\begin{definition}\label{def:fundamental_ideal}
Рассмотрим отображение $\overline\dim\colon W(k)\to\mathbb Z/2\mathbb Z$, индуцированное
размерностью. Ядро этого отображения называется {\bf фундаментальным идеалом} кольца $W(k)$
и обозначается через $IF$.
\end{definition}


%!!! нарисовать коммутативную диаграмму!

\begin{lemma}
%\renewcommand{\theenumi}{\Roman{enumi}}
\begin{enumerate}
\item Идеал $IF$ аддитивно порождается формами вида $\la 1,-a\ra$, $a\in k^*$.
\item $n$-ая степень этого идеала $I^nF:=(IF)^n$ аддитивно порождается тензорными произведениями
$n$ бинарных форм:
$$
\la 1,-a_1\ra\otimes\dots\otimes\la 1,-a_n\ra=:\lla a_1,\dots,a_n\rra.
$$
\end{enumerate}
\end{lemma}
\begin{proof}
Очевидно, что $IF$ порождается формами вида $\la a,b\ra$, $a,b\in k^*$. При этом
$\la a,b\ra\sim\la 1,a\ra\perp -\la 1,-b\ra$. Второе утверждение немедленно следует из первого.
\end{proof}

% 22.03.2010

\begin{definition}
\begin{enumerate}
\item Форма вида $\lla a_1,\dots,a_n\rra$ для $a_1,\dots,a_n\in k^*$
называется {\bf $n$-формой Пфистера} и имеет размерность $2^n$.
Форма называется {\bf формой Пфистера}, если она является $n$-формой Пфистера для некоторого $n$.
\item Если $\ph$~--- форма Пфистера, то $\ph$ представляет 1, поэтому $\ph\cong\la 1\ra\perp-\ph'$
для некоторой формы $\ph'$. Такая форма $\ph'$ называется {\bf чистой формой}, ассоциированной с $\ph$.
\item Обозначим через $P_n(k)$ множество классов изометрий $n$-форм Пфистера; $P(k)=\bigcup_nP_n(k)$;
$GP_n(k)=\{[q]\in W(F)\mid\exists a\in k^*, aq\in P_n(k)\}$; $GP(k)=\bigcup_nGP_n(k)$.
\end{enumerate}
\end{definition}

\begin{lemma}\label{lemma:pfister_identity}
Пусть $a_1,\dots,a_n,b_n\in F^*$. Тогда
$$
\lla a_1,\dots,a_{n-1},a_nb_n\rra\perp\lla a_1,\dots,a_{n-1},a_n,b_n\rra\sim\lla a_1,\dots,a_n\rra\perp\lla a_1,\dots,a_{n-1},b_n\rra.
$$
\end{lemma}
\begin{proof}
Докажем сначала это для $n=1$:
$$
\begin{aligned}
\lla a_1b_1\rra\perp\lla a_1,b_1\rra&\cong\la 1,-a_1b_1,1,-a_1,-b_1,a_1b_1\ra\\
&\sim\la 1,1,-a_1,-b_1\ra\\
&\cong\lla a_1\rra\perp\lla b_1\rra.
\end{aligned}
$$
Остается домножить обе части этого соотношения на $\lla a_1,\dots,a_{n-1}\rra$.
\end{proof}

Пусть $\tilde{IF}$~--- ядро отображения $\dim\colon\tilde W(k)\to\mathbb Z$ и $\tilde{I^nF}=(\tilde{IF})^n$.
Отображение $\tilde W(k)\to W(k)$ индуцирует гомоморфизмы $\tilde{I^nF}\to I^nF$.
\begin{lemma}
Эти гомоморфизмы биективны для всех $n\geq 1$.
\end{lemma}
\begin{proof}
Достаточно рассмотреть случай $n=1$. Для доказательства сюръективности заметим, что $\la 1\ra-\la a\ra$
переходит в $\la 1,-a\ra$. Инъективность: пусть $q,q'$~--- формы одной размерности такие, что $q-q'$
переходит в 0 в кольце $W(k)$. Тогда $q\perp -q'\sim 0$, откуда $q\cong q'$.
\end{proof}

Если $n=1$, форма $\ph=\la 1,-a\ra$ является нормой квадратичного расширения $A=k(\sqrt{a})$ поля $k$.
Если $n=2$, $\ph=\lla a,b\rra$ есть приведенная норма алгебры кватернионов $A=\quaternion{a}{b}{k}$.
Если $n=3$, $\ph=\lla a,b,c\rra$ есть норма неассоциативной алгебры октонионов $A$, определяемой $a,b,c$.
В каждом из этих случаев выполняется тождество $\ph(x\cdot y)=\ph(x)\ph(y)$ для всех $x,y\in A$.
В частности, если $\ph(x)\neq 0$, то $\ph\cong\ph(x)\ph$. Иными словами, $D(\ph)=G(\ph)$.
Если $n\geq 4$, $\ph=\lla a_1,\dots,a_n\rra$ соответствует алгебре Кэли--Диксона $A$, определяемой
$a_1,\dots,a_n$, но не обязательно $\ph(x\cdot y)=\ph(x)\ph(y)$.

\begin{theorem}[Пфистер]\label{theorem:pfister_similar}
Если $\ph$~--- форма Пфистера, $\ph(x)\neq 0$, то $\ph\cong\ph(x)\ph$. Иными словами, $D(\ph)=G(\ph)$
\end{theorem}

\begin{lemma}\label{lemma:pfister_equiv}
Если $a,b,t\in k^*$, то $\lla a,b\rra\cong\lla -ab,a+b\rra$, $\lla a,b\rra\cong\lla a,(t^2-a)b\rra$.
\end{lemma}
\begin{proof}
Заметим, что $\la -a,-b\ra\cong\la -a-b,ab(-a-b)\ra$, откуда
$$
\lla a,b\rra=\la 1,-a,-b,ab\ra\cong\la 1,ab,-a-b,ab(-a-b)\ra=\lla -ab,a+b\rra.
$$
С другой стороны, для 1-форм Пфистера теорема~\ref{theorem:pfister_similar} уже доказана, так что
$\la 1,-a\ra\cong(t^2-a)\la 1,-a\ra$, откуда $\la -b,ab\ra\cong\la -(t^2-a)b,(t^2-a)ab\ra$
и $\lla a,b\rra=\la 1,-a,-b,ab\ra\cong\la 1,-a,-(t^2-a)b,(t^2-a)ab\ra=\lla a,(t^2-a)b\rra$.
\end{proof}

\begin{proposition}\label{prop:pfister_pure_value}
Пусть $\ph=\lla a_1,\dots,a_n\rra$~--- $n$-форма Пфистера и $b\in D(\ph')$, где $\ph'$~--- чистая форма,
ассоциированная с $\ph$. Тогда найдутся $b_2,\dots,b_n\in k^*$ такие, что
$\ph\cong\lla b,b_2,\dots,b_n\rra$.
\end{proposition}
\begin{proof}
Доказываем индукцией по $n$. Если $n=1$, то $b=a_1c^2$ для некоторого $c\in k^*$, и все очевидно.
Предположим, что $n>1$ и обозначим $\tau=\lla a_1,\dots,a_{n-1}\rra\cong\la 1\ra\perp -\tau'$,
тогда $\ph'=\tau'\perp a_n\tau$.
Запишем $b=x+a_ny$ для $x\in D(\tau')\cup\{0\}$, $y\in D(\tau)\cup\{0\}$. Рассмотрим несколько случаев:
\begin{enumerate}
\item Если $y=0$, то $x\neq 0$ и $b\in D(\tau')$, откуда по предположению индукции $\tau\cong\lla b,b_2,\dots,b_{n-1}\rra$,
поэтому $\ph\cong\lla b,b_2,\dots,b_{n-1},a_n\rra$.
\item Если $y\neq 0$, мы покажем, что $\ph\cong\lla a_1,\dots,a_{n-1},a_ny\rra$. Запишем $y=t^2-y_0$ для $y_0\in D(\tau')\cup\{0\}$.
\begin{enumerate}
\item Если $y_0=0$, $y=t^2$ и все очевидно.
\item Если $y_0\in D(\tau')$, то по предположению индукции имеем
$\tau\cong\lla y_0,c_1,\dots,c_{n-1}\rra$, поэтому
$$
\begin{aligned}
\ph&\cong\lla y_0,c_2,\dots,c_{n-1},a_n\rra\\
&\cong\lla y_0,c_2,\dots,c_{n-1},(t^2-y_0)a_n\rra\\
&=\lla y_0,c_2,\dots,c_{n-1},a_ny\rra\\
&\cong\lla a_1,\dots,a_{n-1},a_ny\rra,
\end{aligned}
$$
что и требовалось.
\end{enumerate}
Теперь если $x=0$, то $a_ny=b$ и все в порядке. Если же $x\in D(\tau')$, то $\tau=\lla x,d_2,\dots,d_{n-1}\rra$,
откуда
$$\ph\cong\lla x,d_2,\dots,d_{n-1},a_ny\rra
\cong\lla x+a_ny,d_2,\dots,d_{n-1},-xa_ny\rra
\cong\lla b,\dots\rra.$$
\end{enumerate}
\end{proof}

\begin{corollary}
Изотропная форма Пфистера гиперболична.
\end{corollary}
\begin{proof}
Если $\ph$ изотропна, то $1\in D(\ph')$
\end{proof}
\begin{proof}[Доказательство теоремы~\ref{theorem:pfister_similar}]
Запишем $\ph(x)=t^2-a$ для $a\in D(\ph')\cup\{0\}$. Если $a=0$, утверждение очевидно.
Если $a\in D(\ph')$, то $\ph\cong\lla a\rra\otimes\tau$ для некоторой формы Пфистера $\tau$.
Тогда $\ph(x)\ph=(t^2-a)\lla a\rra\tau\cong\lla a\rra\tau\cong\ph$, поскольку $\lla a\rra$
мультипликативна.
\end{proof}

\begin{corollary}
Две пропорциональные формы Пфистера изометричны.
\end{corollary}
\begin{proof}
Действительно, если $\ph,\ph'$~--- две формы Пфистера, $a\in k^*$ и $a\ph\cong\ph'$, то
из $1\in D(\ph)$ следует $a\in D(\ph')$, откуда $\ph'\cong a\ph'$ по теореме~\ref{theorem:pfister_similar}.
\end{proof}

\begin{corollary}\label{corr:pfister_isotrop}
Пусть $\ph\in P(k)$. Тогда, для всякого $a\in D(\ph)$ и $b\in k^*$ имеем $\ph\otimes\lla a\rra\sim 0$ и
$\ph\otimes\lla ab\rra\cong\ph\otimes\lla b\rra$.
\end{corollary}
\begin{proof}
Первая часть немедленно следует из теоремы~\ref{theorem:pfister_similar}; вторая~--- из первой и
леммы~\ref{lemma:pfister_identity}.
\end{proof}

\begin{corollary}
Пусть $q$~--- квадратичная форма размерности $>1$ над $k$ и $\ph\in P_n(k)$. Предположим, что $q\otimes\ph$ изотропна. Тогда
\begin{enumerate}
\item Найдется изотропная форма $q'$ такая, что $q\otimes\ph\cong q'\otimes\ph$.
\item Анизотропная часть $q\otimes\ph$ имеет вид $\rho\otimes\ph$ для некоторой формы $\rho$.
\item Индекс Витта формы $q\otimes\ph$ делится на $2^n$.
\end{enumerate}
\end{corollary}
\begin{proof}
\begin{enumerate}
\item
Если форма $\ph$ изотропна, то она гиперболична и все очевидно. Пусть $\ph$ анизотропна. Запишем
$q\cong\la a_1,\dots,a_n\ra$. По предположению существуют $b_1,\dots,b_n\in D(\ph)\cup\{0\}$ такие,
что $a_1b_1+\dots+a_nb_n=0$ и не все $b_i$ равны нулю. Не умаляя общности, можно считать, что
$b_1,\dots,b_r\in D(\ph)$ и $b_{r+1}=\dots=b_n=0$. Положим $q'=\la a_1b_1,\dots,a_rb_r,a_{r+1},\dots,a_n\ra$.
Тогда $q'$ изотропна и $q'\otimes\ph\cong a_1b_1\ph\perp\dots\perp a_rb_r\ph\perp a_{r+1}\ph\perp\dots\perp a_n\ph
\cong a_1\ph\perp\dots\perp a_r\ph\perp a_{r+1}\ph\perp\dots\perp a_n\ph\cong q\otimes\ph$.
\item Если $q'$ такая, как в предыдущем абзаце и $m=i(q')$ максимален, то
$q\otimes\ph\cong (m\times\mathbb H\perp\rho)\otimes\ph\sim\rho\otimes\ph$ и $\rho\otimes\ph$
анизотропна по предыдущему пункту.
\item следует из предыдущего.
\end{enumerate}
\end{proof}

\subsection{Суммы квадратов и $s$-инвариант}

\begin{definition}
Пусть $k$~--- поле; $s$-инвариантом $k$ называется наименьшее целое $s(k)$ такое, что $-1$
является суммой $s(k)$ квадратов в $k$. Если такого не существует, полагаем $s(k)=+\infty$.
\end{definition}
\begin{theorem}[Артин--Шрайер]
$s(k)=+\infty$ тогда и только тогда, когда $k$ можно снабдить структурой упорядоченного поля. В этом случае говорят,
что $k$~--- {\bf формально вещественное поле}.
\end{theorem}

\begin{theorem}
Если $s(k)<+\infty$, то это степень двойки.
\end{theorem}
\begin{proof}
Положим $s=s(k)$, пусть $n$~--- целое число такое, что $2^n\leq s<2^{n+1}$. Положим $\ph=\lla -1\rra^{\otimes n}$.
Из определения $s$ следует, что найдутся $x,y$ такие, что $y\neq 0$ и $\ph(x)=-\ph(y)$. При этом $\ph(y)\neq 0$
(иначе $s<2^n$). Значит, $-1=\ph(x)/\ph(y)\in D(\ph)$ по теореме~\ref{theorem:pfister_similar}, откуда $s\leq 2^n$.
\end{proof}


% 29.03.2010

%\bigskip\hrule\bigskip


\begin{definition}
Если $A$~--- абелева группа, {\bf экспонентой} $A$ называется наименьшее целое число $m>0$ такое, что $mA=0$
(или $+\infty$, если такого не существует).
\end{definition}

\begin{proposition}
\begin{enumerate}
\item Экспонента $W(k)$ равна $2s(k)$.
\item Если $s(k)<+\infty$, то всякий элемент $IF$ является нильпотентом. В частности, $W(k)$~--- локальное кольцо
с максимальным идеалом $IF$.
\end{enumerate}
\end{proposition}
\begin{proof}
\begin{enumerate}
\item
Экспонента аддитивной группы кольца равна порядку единицы. Обозначим $s=s(k)$. Это степень двойки, поэтому
достаточно показать, что $s\times\la 1\ra\not\sim 0$ и $2s\times\la 1\ra\sim 0$. Первое следует из определения $s$;
для доказательства второго заметим, что $(s+1)\times\la 1\ra$ является изотропной подформой формы
Пфистера $2s\times\la 1\ra$, которая гиперболична.
\item Для всякой формы $q=\la a_1,\dots,a_n\ra$ размерности $n$ имеем
$q\otimes q\cong n\times\la 1\ra\perp\perp_{i\neq j}\la a_ia_j\ra\cong n\times\la 1\ra\perp\ph\perp\ph$,
где $\ph=\perp_{i<j}\la a_ia_j\ra$. Если $q\in IF$, то $q^2\in 2W(k)$, поэтому
$q^{2r}\in 2^rW(k)$ для всех $r>1$.
\end{enumerate}
\end{proof}


\subsection{Связанные формы Пфистера}

\begin{definition}
Пусть $\ph_1,\ph_2$~--- две формы Пфистера. Будем говорить, что $\ph_1$ и $\ph_2$
являются {\bf $r$-связанными}, если существует $r$-форма Пфистера~$\tau$ и формы
Пфистера $\psi_1$ и $\psi_2$ такие, что $\ph_1\cong\tau\otimes\psi_1$
и $\ph_2\cong\tau\otimes\psi_2$. Формы $\ph_1$ и $\ph_2$ называются
{\bf связанными}, если они являются $(n-1)$-связанными $n$-формами Пфистера.
\end{definition}
\begin{theorem}
Пусть $\ph_1,\ph_2$~--- две анизотропные формы Пфистера и $a_1,a_2\in k^*$.
Тогда $i(a_1\ph_1\perp a_1\ph_2)=0$ или $2^r$, где $r$~--- наибольшее
целое число, для которого $\ph_1$ и $\ph_2$ являются $r$-связанными.
\end{theorem}
Для доказательства теоремы нам потребуется некоторое усиление
предложения~\ref{prop:pfister_pure_value}.
\begin{proposition}\label{prop:pfister_pure_value_strong}
Пусть $\ph\in P_r(k)$, $\psi\in P_s(k)$~--- две формы Пфистера, $\psi'$~--- чистая
форма, ассоциированная с $\psi$. Если $a\in D(\psi'\otimes\ph)$, то существует
$\tau\in P(k)$ такая, что $\psi\otimes\ph\cong\lla a\rra\otimes\tau\otimes\ph$.
\end{proposition}
\begin{proof}
Индукция по $s$. Если $s=1$, то $\psi\cong\lla b\rra$ и $a\in D(b\ph)$,
откуда $ab\in D(\ph)$. По следствию~\ref{corr:pfister_isotrop} имеем
$\lla a\rra\otimes\ph\cong\lla b\rra\otimes\ph$.
Пусть теперь $s>1$, $\psi\cong\lla b\rra\otimes\psi_1$, $\psi'_1$~--- чистая
форма, ассоциированная с $\psi_1$. Тогда
$\psi'\otimes\ph\cong b\psi_1\otimes\ph\perp\psi'_1\otimes\ph$.
Запишем $a=bx+y$, где $x\in D(\psi_1\otimes\ph)\cup\{0\}$, $y\in D(\psi'_1\otimes\ph)\cup\{0\}$.
Предположим сначала, что $x,y\neq 0$. Тогда
$\psi\otimes\ph\cong\lla b\rra\otimes\psi_1\otimes\ph\cong\lla bx\rra\otimes\psi_1\otimes\ph$
по следствию~\ref{corr:pfister_isotrop}.
Кроме того, по предположению индукции, существует форма Пфистера $\tau_1\in P_{s-1}(k)$ такая, что
$\psi_1\otimes\ph\cong\lla y\rra\otimes\tau_1\otimes\ph$. Теперь по лемме~\ref{lemma:pfister_equiv} имеем
$\psi\otimes\ph\cong\lla bx,y\rra\otimes\tau\otimes\ph\cong\lla a,-bxy\rra\otimes\tau_1\otimes\ph$.
Если же $y=0$ или $x=0$, достаточно только половины из этих рассуждений.
\end{proof}
\begin{proof}[Доказательство теоремы]
Пусть $\tau\in P_r(k)$, $\psi_1,\psi_2\in P(k)$ таковы, что $\ph_1\cong\tau\otimes\psi_1$, $\ph_2\cong\tau\otimes\psi_2$
и $r$ максимально. Если форма $a_1\ph_1\perp a_2\ph_2$ анизотропна, доказывать нечего. Если же она изотропна, то
найдется $b\in D(a_1\ph_1)\cap D(-a_2\ph_2)$, откуда $a_1b\in D(\ph_1)$ и $-a_2b\in D(\ph_2)$. Тогда
$a_1\ph_1\cong b\ph_1$ и $a_2\ph_2\cong-b\ph_2$. Теперь не умаляя общности можно считать, что $a_1=1$, $a_2=-1$.
Имеем $\ph_1\perp -\ph_2\sim\tau\otimes(\psi'_2\perp -\psi'_1$, где $\psi'_1$ и $\psi'_2$~--- чистые формы,
ассоциированные с $\psi_1$ и $\psi_2$. При этом $\dim(\ph_1\perp -\ph_2)-\dim(\tau\otimes(\psi'_2\perp -\psi'_1))=2^{r+1}$.
Осталось показать, что форма $\tau\otimes(\psi'_2\perp -\psi'_1)$ анизотропна. Предположим противное.
Тогда $a\in D(\tau\otimes\psi'_1)\cap D(\tau\otimes\psi'_2)$. Но тогда из предложения~\ref{prop:pfister_pure_value_strong}
следует, что $\ph_1$ и $\ph_2$ на самом деле $(r+1)$-связанные, что противоречит выбору $r$.
\end{proof}

\subsection{Мультипликативные формы}

\begin{definition}
Пусть $V$~--- конечномерное векторное пространство над $k$. Мы будем обозначать через $k(V)$
поле частных кольца $\bigoplus_{n\geq 0}S^n(V)$, где $S^n(V)$~--- $n$-ая симметрическая степень $V$.
После выбора базиса $(e_1,\dots,e_n)$ в $V$ поле $k(V)$ отождествляется с полем рациональных
функций от переменных $(e_1,\dots,e_n)$. С точки зрения алгебраической геометрии $k(V)$ является
полем функций аффинного многообразия $\mathbb V$ такого, что $\mathbb V(k)=V^*$~--- пространство,
двойственное к $V$. В частности, если $(V,q)$~--- квадратичное пространство, то $q$ можно
считать элементом $S^2(V^*)$, то есть, элементом $k(V^*)$. Если $(T_1,\dots,T_n)$~--- базис,
двойственный к $(e_1,\dots,e_n)$, то $k(V^*)\cong k(T_1,\dots,T_n)$. Очевидно,
что $q=q(T_1,\dots,T_n)\in D(q_{k(V^*)})$.
\end{definition}

\begin{definition}
Квадратичная форма $\ph$ на пространстве $V$ называется {\bf мультипликативной}, если для
$a=(\ph,0)\in K^*$ и $b=(0,\ph)\in K^*$ имеем $ab\in D(\ph_K)$, где $K=k(V^*\times V^*)$.
Пусть $(T_1,\dots,T_N)$~--- базис $V^*$ и $K=k(U_1,\dots,U_n,V_1,\dots,V_n)$,
где $U_i=(T_i,0)$ и $V_i=(0,T_i)$. Тогда условие мультипликативности можно переформулировать
так: найдутся $f_1,\dots,f_n\in K$ такие, что $\ph(U_1,\dots,U_n)\ph(V_1,\dots,V_n)=\ph(f_1,\dots,f_n)$.
\end{definition}

\begin{theorem}[Классификация анизотропных мультипликативных форм]
Пусть $\ph$~--- анизотропная квадратичная форма над $k$. Следующие условия эквивалентны:
\begin{enumerate}
\item $\ph$ мультипликативна.
\item Для всякого расширения $K/k$ множество $D(\ph_K)$ является подгруппой в $K^*$.
\item Для всякого чисто трансцендентного расширения $K/k$ множество $D(\ph_K)$ является подгруппой в $K^*$.
\item $\ph$ является формой Пфистера.
\end{enumerate}
\end{theorem}

% 5.04.2010

\begin{proof}
$(4)\Rightarrow(2)$~--- из теоремы~\ref{theorem:pfister_similar}, $(2)\Rightarrow(1)$ и $(2)\Rightarrow(3)$~--- очевидно,
$(1)\Rightarrow(2)$~--- из принципа подстановки (применительно к $K$; заметим, что если $q$ мультипликативна, то она остается
мультипликативной после любого расширения $k$). Остается доказать $(3)\Rightarrow(4)$. Пусть $n=\dim(q)$ и $m$~--- наибольшее
целое, для которого $q$ содержит некоторую подформу, изометричную $m$-форме Пфистера. Покажем, что $n=2^m$. Предположим
противное: $n>2^m$, $\ph\leq q$ и $\ph\in P_m(k)$. Запишем $q\cong\ph\perp q'$. Пусть $K=k(V^*\times V^*)$, где $V$~--- пространство,
на котором задана форма $\ph$. По $(3)\Rightarrow(1)$, примененному к $\ph$, есть тождество
$\ph(U)\ph(V)=\ph(f)$, где $f\in K\otimes_kV$. Пусть $a\in D(q')$. Над $K$ имеет место
$$
0\neq\ph(U)+a\ph(V)=\frac{\ph(f)}{\ph(V)}+a\ph(V)=\ph(V)\left(\ph\left(\frac{f}{\ph(V)}\right)+a\right).
$$
Оба множителя справа лежат в $D(q_K)$. Из мультипликативности $q$ следует, что $\ph(U)+a\ph(V)\in D(q_K)$.
Отсюда по теореме о подформе $q\geq\ph\perp a\ph\in P_{m+1}(k)$, что противоречит максимальности $m$.
\end{proof}

Таким образом, если $n$~--- степень двойки, то имеется тождество
$$
(x_1^2+\dots+x_n^2)(y_1^2+\dots+y_n^2)=z_1^2+\dots+z_n^2,
$$
где $z_1,\dots,z_n$~--- рациональные функции от $x_1,\dots,x_n,y_1,\dots,y_n$. На самом деле,
можно доказать, что существуют такие функции $z_i$, линейные по $y$, то есть
$z_i=\sum_j t_{ij}(x_1,\dots,x_n)y_j$, где $t_{ij}\in k(x_1,\dots,x_n)$.

\begin{definition}
Квадратичная форма $\ph$ называется {\bf round-формой}, если $D_k(\ph)=G_k(\ph)$.
\end{definition}

\begin{definition}
Обозначим через $W_t(k)$ {\bf подгруппу кручения} группы $W(k)$:
$W_t(k)=\{w\in W(k)\mid l\times w=0\text{ для некоторого }l\in\mathbb N\}$.
Для $w\in W_t(k)$ наименьшее $l$ такое, что $l\times w=0$, называется
{\bf порядком} $w$.
\end{definition}

\begin{theorem}
$W_t(k)$ является 2-группой, то есть порядок любого элемента $w\in W_t(k)$
является степенью двойки.
\end{theorem}
\begin{proof}
Пусть $w'\in W_t(k)$ имеет порядок $l=2^rs$, где $s$ нечетно и $s>1$.
Тогда $w=2^rw'$ имеет порядок $s$. Выберем анизотропную квадратичную форму
$\ph=\la a_1,\dots,a_m\ra$,
представителем которой является $w$. Тогда $s$~--- наименьшее положительное
число, для которого $s\times\ph\sim 0$.

Возьмем теперь любую степень двойки $n$, большую $m$, и рассмотрим форму
$\psi=\la 1,-\sum_1^nx_i^2\ra$ над $k(x)=k(x_1,\dots,x_n)$, где $x_i$~---
набор переменных над $k$. Нетрудно видеть, что $n\times\psi$~---
изотропная форма Пфистера, поэтому $n\times\psi\sim 0$ над $k(x)$.
Из $s\times\ph\sim 0$ и $n\times\psi\sim 0$ получаем, что
$s\times(\ph\otimes\psi)\sim 0$ и $n\times(\ph\otimes\psi)\sim 0$,
откуда $\ph\otimes\psi\sim 0$, поскольку $s$ и $n$ взаимно просты.
Это означает, что $\ph\cong(\sum_1^nx_i^2)\ph$ над $k(x)$.
В частности, $\ph$ представляет элемент $a_1\sum_1^nx_i^2$ над $k(x)$.
Но $a_1\sum_1^nx_i^2$~--- это общий элемент квадратичной формы $n\times\la a_1\ra$.
Поскольку $\ph$ анизотропна над $k$, из теоремы~\ref{theorem:subform} о подформе теперь следует, что
$\ph$ содержит $n\times\la a_1\ra$, поэтому $m=\dim\ph\geq n$~--- противоречие.
\end{proof}

\begin{remark}
С помощью такого же типа рассуждений (со ссылкой на теорему о подформе) нетрудно
доказать (упражнение!), что, скажем, выражение $x^2+y^2+z^2+t^2$ не может быть представлено в виде
суммы {\it трех} квадратов рациональных функций от переменных $x,y,z,t$.
\end{remark}

\begin{theorem}
Пусть $\ph\leq\psi$~--- две формы Пфистера. Тогда существует форма Пфистера $\tau$ такая, что
$\psi=\ph\otimes\tau$.
\end{theorem}
\begin{proof}
Доказывается аналогично; индукция по $\dim\psi-\dim\ph$. Пусть $\psi=\ph\perp q$ и $a\in D(q)$;
нетрудно показать, что $\ph\otimes\la 1,a\ra\leq\psi$.
\end{proof}

\begin{definition}
Пусть $(V,q)$~--- квадратичное пространство размерности $n$ и $Q$~--- матрица формы $q$ в некотором базисе.
{\bf Дискриминантом} формы $q$ называется элемент $d(q)=(-1)^{\frac{n(n-1)}{2}}\det Q\in k^*/(k^*)^2$.
Он не зависит от выбора базиса в $V$.
\end{definition}

\begin{proposition}
Отображение $e_0\colon W(k)\to\mathbb Z/2$, $e_0(\ph)=\dim(\ph)\mod 2$, является сюръективным
гомоморфизмом колец. Ядро этого гомоморфизма~--- фундаментальный идеал $IF$, поэтому
$W(k)/IF\cong\mathbb Z/2$.
\end{proposition}
\begin{proof}
См.~определение~\ref{def:fundamental_ideal}.
\end{proof}

Множитель $(-1)^{\frac{n(n-1)}{2}}$ позволяет дискриминанту быть корректно определенным на
кольце Витта~--- дискриминант не меняется при замене формы на эквивалентную.
Заметим, что при этом дискриминант не является гомоморфизмом: вообще говоря, неверно, что
$d(q_1\perp q_2)=d(q_1)d(q_2)$. Но если ограничиться рассмотрением форм четной размерности
(то есть представителей элементов из $IF$), то оказывается, что дискриминант является
гомоморфизмом,
поскольку для формы $\ph$ четной размерности $d(\ph)=(-1)^{\frac{\dim(\ph)}{2}}\det(\ph)$
Чуть позднее (теорема~\ref{thm:discriminant_iso})
мы докажем, что дискриминант отождествляет $IF/I^2F$ с $k^*/(k^*)^2$.

\begin{definition}
Степени фундаментального идеала определяют фильтрацию кольца Витта.
Пусть $\overline{I^nF}=I^nF/I^{n+1}F$ (при этом $I^0F=W(k)$). Построим
абелеву группу $gr(W)=\overline{I^0F}\oplus\overline{I^1F}\oplus\overline{I^2F}\oplus\dots$
и введем на этой группе операцию умножения, индуцированную умножением в кольце Витта $W(k)$:
для $\overline{x}\in\overline{I^mF}$, $\overline{y}\in\overline{I^nF}$ элемент
$\overline{x}\cdot\overline{y}=\overline{xy}\in\overline{I^{m+n}F}$ корректно определен.
Полученное кольцо называется {\bf градуированным кольцом Витта} поля~$k$.
\end{definition}

Оказывается, что $\overline{I^2F}$ отождествляется с 2-кручением группы
Брауэра~--- классического объекта. В ближайшее время мы построим по форме $q$ центральную простую алгебру
(Клиффорда), и сопоставление форме класса этой алгебры в группе Брауэра поля $k$
окажется гомоморфизмом $e_2\colon I^2F\to Br(k)$, который превратится в
изоморфизм $\overline{e_2}\colon \overline{I^2F}\to {}_2Br(k)$.

\section{$K$-теория Милнора}

% 12.04.2010

\subsection{Элементарные инварианты}

\begin{proposition}\label{prop:disc_properties}
\begin{enumerate}
\item Если $q=\la a_1,\dots,a_n\ra$, то $d_q=(-1)^{\frac{n(n-1)}{2}}a_1\dots a_n$.
\item $d(q\perp q')=d_qd_{q'}(-1)^{nn'}$, где $n'=\dim q'$.
\item $d(q\perp\mathbb H)=d_q$.
\item $d(q\otimes q')=(d_q)^{n'}(d_{q'})^n$.
\end{enumerate}
\end{proposition}
\begin{proof}
Легко.
\end{proof}

\begin{theorem}\label{thm:discriminant_iso}
Инвариант $d$ индуцирует изоморфизм $IF/I^2F\to k^*/(k^*)^2$.
\end{theorem}
\begin{proof}
Из предыдущего предложения следует, что $d$ является гомоморфизмом
и $d|_{I^2F}=1$. Значит, $d$ индуцирует гомоморфизм
$\overline{d}\colon IF/I^2F\to k^*/(k^*)^2$. Он сюръективен,
поскольку $d(\la 1,-a\ra)=a$ для $a\in k^*/(k^*)^2$. Для доказательства
инъективности предположим, что $q\in IF$ и $d(q)=1$. Ведем индукцию
по $2n=\dim q$. Если $n=1$, то $q=\la a_1,a_2\ra$ и $a_2=-a_1$ по модулю $(k^*)^2$,
откуда $\ph\cong a_1\la 1,-1\ra$ и $\tilde\ph=0$ в $W(k)$. Если $n>1$, то
$\ph=\la a_1,a_2,a_3\ra\perp\la a_4,\dots,a_{2n}\ra\sim\la a_1,a_2,a_3,a_1a_2a_3\ra\perp
\la -a_1a_2a_3,a_4,\dots,a_{2n}\ra$. Заметим, что $\la a_1,a_2,a_3,a_1a_2a_3\ra\cong
\la a_1,a_2\ra\otimes\la 1,a_1a_3\ra\in I^2F$, размерность формы
$\la -a_1a_2a_3,a_4,\dots,a_{2n}\ra$ равна $2(n-1)$, а дискриминант равен 1.
Значит, она лежит в $I^2F$ по предположению индукции, откуда $\ph\in I^2F$.
\end{proof}

\begin{corollary}
\begin{enumerate}
\item Если размерность формы $q$ нечетна, то $q\equiv\la d(q)\ra\pmod{I^2F}$.
\item Если размерность формы $q$ четна, то $q\equiv\la 1,-d(q)\ra\pmod{I^2F}$.
\end{enumerate}
\end{corollary}
\begin{proof}
Очевидно.
\end{proof}

Можно явно описать расширение $\mathbb Z/2$ с помощью $k^*/(k^*)^2$, определенное $W(k)/I^2F$.
Обозначим через $Q(k)$ множество $\mathbb Z/2\otimes k^*/(k^*)^2$, снабженное следующей операцией:
$$
(a,u)+(b,v)=(a+b,(-1)^{ab}uv).
$$

\begin{proposition}\label{prop:group_Q}
Отображение $q\mapsto (\overline{\dim}(q),d(q))$ индуцирует изоморфизм
$$
W(k)/I^2F\to Q(k).
$$
\end{proposition}
\begin{proof}
Достаточно проверить, что это гомоморфизм; это напрямую следует из
предложения~\ref{prop:disc_properties}.
\end{proof}

\begin{definition}
Пусть $(V,q)$~--- квадратичное пространство над $k$. {\bf Алгеброй Клиффорда} $C(q)$
формы $q$ называется фактор-алгебра тензорной алгебры $T(V)$ пространства $V$ по
двустороннему идеалу, порожденному элементами вида $v\otimes v-q(v)1$, $v\in V$.
\end{definition}

Если $q=\la a_1,\dots,a_n\ra$ в ортогональном базисе $(e_1,\dots,e_n)$ пространства $V$,
то $C(q)$ можно описать как алгебру, порожденную элементами $e_i$ с соотношениями
$e_i^2=a_i$ и $e_ie_j+e_je_i=0$ для $i\neq j$.

\begin{proposition}[универсальное свойство $C(q)$]
Если $A$~--- $k$-алгебра, $f\colon V\to A$~--- гомоморфизм векторных пространств
над $k$ такой, что $f(v)^2=q(v)$ для всех $v\in V$, то $f$ единственным образом
продолжается до гомоморфизма $k$-алгебр $\tilde f\colon C(V)\to A$.
\end{proposition}

Поскольку $T(V)$ является градуированной алгеброй и соотношения в $C(q)$ однородны
по модулю 2, то алгебра $C(q)$ обладает естественной $\mathbb Z/2$-градуировкой.
Будем обозначать через $C_0(q)$ (соответственно $C_1(q)$) ее четную (соответственно нечетную) часть.
Алгебру с $\mathbb Z/2$-градуировкой еще называют {\bf супералгеброй}.

\begin{definition}
Пусть $A$, $B$~--- две супералгебры над $k$. {\bf Градуированным тензорным произведением} $A$ и $B$
называется супералгебра $A\hat\otimes_k B$ такая, что
\par\noindent $\bullet$ $A\hat\otimes_k B$ совпадает с $A\otimes_k B$ как векторное пространство;
\par\noindent $\bullet$ если $(a,a',b,b')\in A^2\times B^2$~--- однородны, то
$$
(a\hat\otimes b)(a'\hat\otimes b')=(-1)^{|a'||b|}aa'\hat\otimes bb';
$$
\par\noindent $\bullet$ если $a\in A$, $b\in B$ однородны степеней $i$, $j$, то $ab$ однороден
степени $i+j$.
\end{definition}

\begin{proposition}
Если $\dim q=n$, то $\dim_k C(q)=2^n$.
\end{proposition}

Если $(V_1,q_1)$, $(V_2,q_2)$~--- два квадратичных пространства над $k$, включения
$V_i\hookrightarrow V_1\otimes V_2\hookrightarrow C(q_1\perp q_2)$ вместе с универсальным
свойство алгебры Клиффорда индуцируют гомоморфизмы алгебр $C(q_i)\to C(q_1\perp q_2)$,
которые являются и гомоморфизмами супералгебр.
В $C(q_1\perp q_2)$ выполнено $v_1v_2=-v_2v_1$ для $(v_1,v_2)\in V_1\times V_2$;
эти гомоморфизмы продолжаются до гомоморфизма супералгебр
$C(q_1)\hat\otimes_kC(q_2)\to C(q_1\perp q_2)$.
\begin{theorem}\label{thm:tensor_product_of_superalgebras}
Этот гомоморфизм является изоморфизмом супералгебр.
\end{theorem}
\begin{proof}
Сюръективность очевидна; инъективность следует из предыдущего предложения и соображений размерности.
\end{proof}

\begin{corollary}
Пусть $q$~--- квадратичная форма, $a\in k^*$ и $q'=\la -a\ra\perp q$.
Тогда $C(aq)$ изоморфна (как алгебра) $C_0(q')$.
\end{corollary}
\begin{proof}
$C(q')\cong C(\la -a\ra)\hat\otimes_k C(q)$, откуда
$$
C(q')\cong C_0(\la -a\ra)\hat\otimes_kC(q)\oplus C_1(\la -a\ra)\hat\otimes_kC(q)=C(q)\oplus zC(q)
$$
(изоморфизмы векторных пространств), где $z$~--- каноническая образующая $C_q(\la-a\ra)$, $z^2=-a$.
Отсюда $C_0(q')\cong C_0(q)\oplus zC_1(q)$.
Остается отождествить последнее слагаемое с $C(aq)$. Но $z$ коммутирует с $C_0(q)$
и антикоммутирует с $C_1(q)$; в частности, $(zv)^2=zvzv=-z^2v^2=aq(v)$ для всех $v\in V$.
Из универсального свойства алгебры Клиффорда теперь следует, что линейное отображение
$V\to C_0(q)\oplus zC_1(q)$, $v\mapsto zv$ продолжается до гомоморфизма алгебр
$C(aq)\to C_0(q)\oplus zC_1(q)$. Очевидно, что этот гомоморфизм сюръективен,
и биективен по соображениям размерности.
\end{proof}

\begin{corollary}
$\dim C_0(q)=\dim C_1(q)=2^{n-1}$.
\end{corollary}

\subsection{Группа Брауэра}

\begin{definition}
Пусть $A$~--- кольцо. $A$-модуль $M$ называется {\bf простым}, если у него нет
подмодулей, кроме $M$ и $0$. $M$ называется {\bf полупростым}, если он удовлетворяет
следующим эквивалентным условиям:
\begin{enumerate}
\item $M$~--- сумма своих простых подмодулей;
\item $M$~--- прямая сумма простых модулей;
\item всякий подмодуль $M$ выделяется прямым слагаемым.
\end{enumerate}
\end{definition}

\begin{definition}
Кольцо $A$ называется {\bf полупростым}, если оно удовлетворяет
следующим эквивалентным условиям:
\begin{enumerate}
\item всякий левый $A$-модуль прост;
\item $A$ прост как левый $A$-модуль;
\item всякий идеал $A$ выделяется прямым слагаемым.
\end{enumerate}
Полупростое кольцо $A$ называется {\bf простым}, если в нем нет двусторонних идеалов,
отличных от $0$ и $A$.
\end{definition}

\begin{theorem}
Всякое полупростое кольцо является прямым произведением конечного числа простых колец.
Оно является простым тогда и только тогда, когда всякий простой модуль над ним
изоморфен ему. Всякое простое кольцо изоморфно алгебре матриц над телом.
\end{theorem}

\begin{definition}
Алгебру над полем $F$ будем называть {\bf простой $F$-алгеброй}, если она конечномерна
над $F$ и является простым кольцом. $F$-алгебра называется {\bf центральной},
если ее центр совпадает с $F$.
\end{definition}
\begin{definition}
Пусть $A$~--- $F$-алгебра, $B$~--- подалгебра $A$. {\bf Централизатор} $B$ в $A$~---
это $B'=\{a\in A\mid ab=ba\;\forall b\in B\}$. $B'$ является подалгеброй в $A$.
\end{definition}
\begin{theorem}
\begin{enumerate}
\item Пусть $K/F$~--- расширение полей. $F$-алгебра $A$ является центральной простой тогда и
только тогда, когда $K$-алгебра $A_K:=K\otimes_F A$ является центральной простой.
\item Пусть $A$~--- центральная простая $F$-алгебра.
\begin{enumerate}
\item Размерность $A$ над $F$ является точным квадратом.
\item Пусть $B$~--- простая $F$-подалгебра $A$, $B'$~--- ее централизатор в $A$. Тогда
\begin{enumerate}
\item $B'$ проста.
\item $[B:F][B':F]=[A:F]$.
\item централизатор $B'$ в $A$ совпадает с $B$.
\item если $B$ центральна, то $B'$ центральна и гомоморфизм $B\otimes_F B'\to A$ является изоморфизмом.
\end{enumerate}
\item Пусть $B$~--- центральная простая $F$-алгебра. Тогда алгебра $A\otimes_F B$ является
центральной простой.
\item Пусть $A^0$~--- алгебра, противоположная к $A$. Тогда существует канонический изоморфизм
$A\otimes_F A^0\cong\End_F(A)$.
\end{enumerate}
\end{enumerate}
\end{theorem}

% 19.04.2010
% группа Брауэра

\begin{definition}
Пусть $F$~--- поле. Две конечномерные центральные простые $F$-алгебры $A,B$ называются {\bf подобными},
если выполняются следующие эквивалентные условия:
\begin{enumerate}
\item Для некоторого тела $D$ с центром $F$ и целых $a,b$ выполняется $A\cong M_a(D)$ и $B\cong M_b(D)$.
\item Найдутся целые $a,b$ такие, что $M_b(A)\cong M_a(B)$.
\item Категории левых $A$-модулей и левых $B$-модулей эквивалентны.
\end{enumerate}
Будем говорить, что $A$ {\bf нейтральна}, если $A$ подобна $F$.
\end{definition}
Из третьего условия видно, что это подобие является отношением эквивалентности.

\begin{theorem}
\begin{enumerate}
\item Совокупность классов подобия центральных простых $F$-алгебр образует множество $Br(F)$.
\item Тензорное произведение снабжает $Br(F)$ структурой группы; нейтральным элементом является
класс нейтральных алгебр; обратный к классу алгебры $A$~--- это класс противоположной алгебры $A^0$;
эта группа коммутативна.
\item Если $K/F$~--- расширение полей, то расширение скаляров индуцирует гомоморфизм $Br(F)\to Br(K)$.
\end{enumerate}
\end{theorem}
\begin{proof}
Очевидно.
\end{proof}

\begin{remark}
Группа $Br(F)$ называется {\bf группой Брауэра} поля $F$.
\end{remark}

\begin{examples}
\begin{enumerate}
\item Если поле $F$ алгебраически замкнуто, то $Br(F)=0$.
\item $Br(\mathbb R)\cong\mathbb Z/2$; эту группу порождает класс кватернионов
Гамильтона $\quaternion{-1}{-1}{\mathbb R}$
\item $Br(\mathbb F_q)=0$.
\end{enumerate}
\end{examples}

Заметим, что $\lla a\rra\perp\lla b\rra\equiv\lla ab\rra\pmod{I^2F}$
(см. лемму~\ref{lemma:pfister_identity}).
Кроме того, $\lla x,1-x\rra\cong\lla 1,\dots\rra\sim 0$.
Рассмотрим абелеву группу, порожденную символами
$\{a_1,\dots,a_n\}$, где $a_i\in k^*$, с соотношениями
\begin{enumerate}
\item $\{\dots,a,\dots\}+\{\dots,b,\dots\}=\{\dots,ab,\dots\}$;
\item $\{\dots,x,\dots,1-x,\dots\}=0$.
\end{enumerate}
Введем на этой группе умножение так:
$$
\{a_1,\dots,a_n\}\cdot\{b_1,\dots,b_m\}=\{a_1,\dots,a_n,b_1,\dots,b_m\}.
$$
Получится градуированное кольцо, которое обозначается через $K_*^M(k)$ и называется
{\bf $K$-теорией Милнора} поля $k$. Эквивалентно,
$$
K_*^M(k)\cong T(k^*)/x\otimes(1-x)
$$
--- фактор-алгебра
тензорной алгебры (над $\mathbb Z$) мультипликативной группы поля $k$.

\begin{definition}
Пусть $A$~--- центральная простая $F$-алгебра. Расширение $K/F$ называется
{\bf нейтрализующим полем}, если $A_K$ нейтральна.
\end{definition}

\begin{theorem}
Пусть $A$~--- центральная простая $F$-алгебра; $E/F$~---- конечное расширение.
$E$ является нейтрализующим полем для $A$ тогда и только тогда, когда существует
центральная простая $F$-алгебра $B$, подобная $A$, такая, что $E$ является максимальной
коммутативной подалгеброй в $B$. Более того, следующие утверждения эквивалентны:
\begin{enumerate}
\item $E$~--- максимальная коммутативная подалгебра в $B$.
\item $E$ совпадает со своим централизатором.
\item $[B:F]=[E:F]^2$.
\end{enumerate}
\end{theorem}

\begin{theorem}[Сколем--Нетер]
Пусть $A$~--- центральная простая $F$-алгебра, $B,C$~--- две простые подалгебры $A$
и $f\colon B\to C$~--- изоморфизм $F$-алгебр. Тогда существует обратимый $a\in A$ такой,
что $f(x)=axa^{-1}$ для всех $x\in B$. В частности, любой автоморфизм алгебры $A$
является внутренним.
\end{theorem}

%!!! кажется, этой теоремы не было!
%\begin{theorem}
%Пусть $D$~--- центральное тело над $F$. Существует максимальное коммутативное
%под-тело в $D$, сепарабельное над $F$.
%\end{theorem}

\begin{definition}
Пусть $A$~--- центральная простая $F$-алгебра. Запишем $A=M_n(D)$ для некоторого тела $D$.
\begin{enumerate}
\item {\bf Степень} $A$~--- это целое число $\sqrt{[A:F]}$.
\item {\bf Индекс} $A$~--- это целое число $\sqrt{[D:F]}$.
\item {\bf Экспонента} $A$~--- это порядок класса алгебры $A$ в $Br(F)$.
\end{enumerate}
\end{definition}

\begin{proposition}
Для всякой центральной простой $F$-алгебры $A$
\begin{enumerate}
\item экспонента $A$ делит ее индекс, а индекс делит ее степень;
\item индекс и экспонента $A$ состоят из одинаковых простых делителей. % этого пункта тоже не было
\end{enumerate}
\end{proposition}

\begin{proposition}
Пусть $A$~--- центральная простая $F$-алгебра; $K/F$~--- расширение полей.
\begin{enumerate}
\item $\ind(A)$ делится на $\ind(A_K)$.
\item Если $[K:F]=n<+\infty$, то $\ind(A)/\ind(A_K)$ является делителем $n$.
\item Если $K/F$~--- чисто трансцендентное расширение, то $\ind(A_K)=\ind(A)$.
\end{enumerate}
\end{proposition}

\begin{lemma}[Альберт]
Пусть $F$~--- поле характеристики не 2, $D$~--- конечномерное центральное тело над $F$ и $a\in F^*\setminus (F^*)^2$.
Пусть $E=F(\sqrt{a})$. Тогда $D_E$ не является телом тогда и только тогда, когда $D$ содержит подполе, изоморфное $E$.
\end{lemma}
\begin{proof}
$D$ содержит подполе, изоморфное $E$ тогда и только тогда, когда $a$ является квадратом в $D$. Тогда
$$
(1\otimes x-\sqrt{a}\otimes 1)(1\otimes x+\sqrt{a}\otimes 1)=0,
$$
поэтому $D_E$ содержит делители нуля
и не может быть телом. Обратно, если $D_E$~--- не тело, то найдутся $s,t,u,v\in D$, не все равные нулю,
такие, что $(1\otimes s + \sqrt{a}\otimes t)(q\otimes u + \sqrt{a}\otimes v)=0$. Тогда $s,u\neq 0$ и,
после домножения слева на $s^{-1}$ и справа на $u^{-1}$, можно считать, что $s=u=1$. Тогда
$0=(1+t\sqrt{a})(1+v\sqrt{a})=1+tvd+(t+v)\sqrt{a}$. Отсюда $v=-t$ и $1-dt^2=0$; теперь видно, что $(t^{-1})^2=a$.
\end{proof}

\begin{proposition}
Пусть $A$~--- центральная простая алгебра над $F$, $K$~--- подполе в $A$ и
$B=K'$. Тогда $A_K$ подобна $B$.
\end{proposition}

\begin{definition}
Пусть $a,b\in F^*$. {\bf Алгеброй кватернионов}, построенной по паре $(a,b)$, называется $F$-алгебра
$\quaternion{a}{b}{F}$ с базисом $(1,i,j,k)$ такая, что $i^2=a, j^2=b, ij=-ji=k$.
\end{definition}

\begin{remark}
Алгебра $\quaternion{a}{b}{F}$ совпадает с алгеброй $C(\la a,b\ra)$ (если забыть про структуру супералгебры на этой
алгебре Клиффорда).
\end{remark}

\begin{theorem}\label{thm:quaternion_splitting}
Пусть $a,b\in F^*$. Следующие условия эквивалентны:
\begin{enumerate}
\item Квадратичная форма $\la 1,-a,-b\ra$ изотропна.
\item Квадратичная форма $\lla a,b\rra$ изотропна.
\item Алгебра $Q=\quaternion{a}{b}{F}$ не является телом.
\item Алгебра $Q$ изоморфна $M_2(F)$.
\end{enumerate}
В частности, $Q$~--- центральная простая алгебра над $F$ степени 2.
\end{theorem}
\begin{proof}
$(1)\Longleftrightarrow (2)$~--- нетрудно. %???
$(2)\Longleftrightarrow (3)$: заметим, что если $x,y,z,t\in F$,
то $(x+yi+zj+tk)(x-yi-zj-tk)=x^2-ay^2-bz^2+abt^2=q(x,y,z,t)$, где $q=\lla a,b\rra$.
Значит, если $q$ изотропна, то в $Q$ есть делители нуля, а если $q$ анизотропна,
то всякий элемент $x+yi+zj+tk\in Q\setminus\{0\}$ обратим; обратный к нему равен
$(x-yi-zj-tk)/q(x,y,z,t)$.

Очевидно, что $(4)\Rightarrow(3)$. Покажем что $(3)\Rightarrow(4)$. Предположим,
что $a=1,$ $b=-1$ и построим изоморфизм $\quaternion{a}{b}{F}\cong M_2(F)$.
Пусть $(E_{ij})_{i,j\in\{0,1\}}$~--- канонический базис $\End_F(F\hat\oplus F)$.
Тогда нужный изоморфизм устанавливается так:
$$
\begin{aligned}
% нарисовать в матрицах
1&\mapsto E_{00}+E_{11}\\
i&\mapsto E_{01}+E_{10}\\
j&\mapsto E_{01}-E_{10}\\
k&\mapsto E_{11}-E_{00}\\
\end{aligned}
$$
Заметим, что $\quaternion{a}{b}{F}\cong\quaternion{as^2}{bt^2}{F}$ для $s,t\in F^*$. Значит, для произвольных
$a,b$ существует расширение $E/F$ такое, что $Q_E\cong\quaternion{1}{-1}{E}$. Значит, и $Q$
является центральной простой и, очевидно, степени 2. Если $Q$ не является телом, то она обязана быть
изоморфна $M_2(F)$.
\end{proof}
Верно и обратное:
\begin{theorem}
Всякая центральная простая $F$-алгебра $A$ степени 2 является алгеброй кватернионов.
\end{theorem}
\begin{proof}
Можно считать, что $A$~--- тело. Пусть $E\subset A$~--- максимальное коммутативное подтело $A$.
Тогда $E=F(\sqrt{a})$ для подходящего $a\in F^*$. Возьмем $i\in E$ такое, что $i^2=a$. Рассмотрим
внутренний автоморфизм $\sigma$ алгебры $A$, определенный $i$: видно, что $\sigma^2=1$. Если
$i$ не централен, то $\sigma\neq 1$; поэтому у $\sigma$ есть собственное число, равное $-1$,
то есть, найдется $j\in A$ такой, что $\sigma(j)=-j$. Стало быть, $ij=-ji$. Легко видеть, что
$j$ не централен, поэтому $j$ порождает максимальное коммутативное подтело $K$ в $A$.
Автоморфизм $\sigma$ переводит $K$ в себя и его ограничение на $F$ тривиально; поэтому
множество неподвижных точек $\sigma|_K$ совпадает с $F$. В частности, $j^2=b\in F$.
Наконец, положим $k=ij$. Если $x+yi+zj+tk=0$~--- какая-то нетривиальная линейная комбинация,
то после сопряжения при помощи $i$, $j$ и $k$ получим соотношения
$$
\begin{aligned}
x+yi-zj-tk&=0\\
x-yi+zj-tk&=0\\
x-yi-zj+tk&=0\\
\end{aligned}
$$
откуда $x=y=z=t=0$.
\end{proof}

\begin{lemma}\label{lemma:relations_quaternion}
Для $a,b\in F^*$ обозначим через $(a,b)$ класс алгебры $\quaternion{a}{b}{F}$ в $Br(F)$. Тогда
$$
\begin{aligned}
(a,b)&=(b,a)\\
(a^2,b)&=0\\
(a,1-a)&=0\quad(a\neq 1)\\
(a,-a)&=0\\
(a,bb')&=(a,b)+(a,b')\\
\end{aligned}
$$
\end{lemma}
\begin{proof}
Первое свойство очевидно. Следующие три получаются из теоремы~\ref{thm:quaternion_splitting}.
Для доказательства последнего построим изоморфизм
$$
\ph\colon\quaternion{a}{b}{F}\times_F\quaternion{a}{b'}{F}\stackrel{\sim}{\longrightarrow} M_2\left(\quaternion{a}{bb'}{F}\right).
$$
Пусть $(1,i,j,k)$, $(1,i',j',k')$, $(1,i'',j'',k'')$~--- соответственно канонические базисы алгебр
$\quaternion{a}{b}{F}$, $\quaternion{a}{b'}{F}$ и $\quaternion{a}{bb'}{F}$. Вот
образы некоторых элементов при нашем изоморфизме:
$$
\begin{aligned}
\ph(i\otimes 1)&=\begin{pmatrix}i''&0\\0&i''\end{pmatrix}\qquad & 
\ph(1\otimes i')&=\begin{pmatrix}-i''&0\\0&i''\end{pmatrix}\\
\ph(j\otimes 1)&=\begin{pmatrix}0&-j''\\-b'^{-1}j''&0\end{pmatrix}\qquad & 
\ph(1\otimes j')&=\begin{pmatrix}0&b'\\1&0\end{pmatrix}\\
\ph(k\otimes 1)&=\begin{pmatrix}0&-k''\\-b'^{-1}k''&0\end{pmatrix}\qquad & 
\ph(1\otimes k')&=\begin{pmatrix}0&-b'i''\\i''&0\end{pmatrix}.\\
\end{aligned}
$$
Необходимо лишь проверить, что $\ph$ является гомоморфизмом алгебр; любой гомоморфизм из одной простой алгебры
в другую является инъективным, и сюръективность вытекает из соображений размерности.
\end{proof}

\begin{lemma}[Альберт]\label{lemma:chain_albert}
Пусть $a,b,c,d\in F^*$ таковы, что $(a,b)=(c,d)$. Тогда существует $e\in F^*$ такой, что $(a,b)=(a,e)=(c,e)=(c,d)$.
\end{lemma}
\begin{proof}
Пусть $D=\quaternion{a}{b}{F}$ и $D_0$~--- векторное $F$-подпространство в $D$, состоящее из элементов следа $0$
(ортогональных к 1 по отношению к приведенной норме); тогда $\dim D_0=3$. Ограничение $q$ на $D_0$ совпадает
с отображением $x\mapsto -x^2$. По предположению, найдутся $\alpha,\beta,\gamma,\delta\in D_0$ такие, что
$$
\begin{array}{c}
\alpha^2=a,\beta^2=b,\alpha\beta+\beta\alpha=0\\
\gamma^2=c,\delta^2=d,\gamma\delta+\delta\gamma=0.
\end{array}
$$
Иными словами, $(\alpha,\beta)$ и $(\gamma,\delta)$~--- пары ортогональных векторов. Но $\dim D_0=3$, поэтому
найдется $\varepsilon\in D_0$, ортогональный к $\alpha$ и $\gamma$. Тогда можно взять $e=\varepsilon^2$.
\end{proof}

% 26.04.2010
% группа Брауэра-Уолла

\subsection{Группа Брауэра--Уолла}

В этом разделем мы определим аналог группы Брауэра для супералгебр. Пусть $A$~--- супералгебра над $F$.

\begin{definition}
Супералгебра $A$ называется {\bf простой}, если она не имеет градуированных двусторонних идеалов, кроме $0$ и $A$.
\end{definition}
\begin{definition}
{\bf (Градуированным) центром} супералгебры $A$ называется $\hat{Z}(A)=Z_0(A)\oplus Z_1(A)$, где
$Z_i(A)=\{a\in A_i\mid\forall x\in A_j, ax=(-1)^{ij}xa,j=0,1\}$. Супералгебра $A$ над $F$ называется
{\bf центральной}, если $\hat{Z}(A)=(F,0)$.
\end{definition}

\begin{examples}\label{examples:superalgebras}
\begin{enumerate}
\item Пусть $A$~--- алгебра над $F$. Определим супералгебру $i(A)$ так: $i(A)_0=A$, $i(A)_1=0$.
Если $A$ центральная простая алгебра, то $i(A)$~--- центральная простая супералгебра.
\item Пусть $V=V_0\hat\oplus V_1$~--- конечномерное векторное суперпространство над $F$. Алгебра
$\End_F(V)$ допускает естественную градуировку, в которой эндоморфизм $u$ является четным,
если $u(V_i)\subset V_i$ и нечетным, если $u(V_i)=V_{i+1}$ для всех $i\in\mathbb Z/2$.
Полученная супералгебра является центральной простой.
\item Пусть $a\in F^*$. Супералгебра $C(\la a\ra)$ изоморфна (как алгебра без градуировки)
$F[t]/(t^2-a)$. Ее градуировка определяется однозначно из условия $|t|=1$. Нетрудно видеть,
что эта супералгебра является центральной простой.
\end{enumerate}
\end{examples}

\begin{proposition}\label{prop:tensor_product_of_central_simple_superalgebras}
Если $A,B$~--- центральные простые $F$-супералгеб\-ры, то $A\hat\otimes_F B$~--- тоже
центральная простая супералгебра.
\end{proposition}

\begin{theorem}\label{thm:clifford_superalgebra_is_central_simple}
Для всякой невырожденной квадратичной формы $q$ супералгебра $C(q)$ над $F$ является центральной простой.
\end{theorem}
\begin{proof}
Приведение формы к диагональному виду с учетом теоремы~\ref{thm:tensor_product_of_superalgebras} и
предложения~\ref{prop:tensor_product_of_central_simple_superalgebras} сводит задачу к случаю $\dim q=1$,
который приведен в примере~\ref{examples:superalgebras} (3).
\end{proof}

\begin{definition}
Две $F$-супералгебры $A$, $B$ называются {\bf подобными}, если существуют два векторных
суперпространства $V$, $W$ над $F$ такие, что $A\hat\otimes\End_F(V)\cong B\hat\otimes\End_F(W)$
(см. пример~\ref{examples:superalgebras} (2)).
Обозначение: $A\sim B$.

Пусть $A$~---супералгебра. {\bf Противоположная супералгебра} $A^*$ определяется так:
как векторное пространство $A^*=\{a^*\mid a\in A\}$, градуировка вводится так, что
$A^*_i=\{a^*\mid |a|=i\}$, а произведение выглядит так: $a^*b^*=(-1)^{|a||b|}(ba)^*$
для однородных $a,b$.
\end{definition}

\begin{theorem}
Отношения подобия является отношением эквивалентности на множестве центральных простых
$F$-супералгебр, совместимым с градуированным тензорным произведением. Полугруппа
классов эквивалентности является коммутативной группой и называется
{\bf группой Брауэра--Уолла} поля $F$ и обозначается через $BW(F)$.
Если $A$~--- центральная простая $F$-супералгебра, класса $\la A\ra\in BW(F)$,
то представителем класса $-\la A\ra$ является супералгебра $A^*$, противоположная к $A$.
\end{theorem}
\begin{proof}
Для проверки коммутативности $BW(F)$ заметим, что если $A$ и $B$~--- супералгебры над $F$,
то имеется изоморфизм $F$-супералгебр $A\hat\otimes B\stackrel{\sim}{\longrightarrow}B\hat\otimes A$,
задаваемый так: $a\hat\otimes b\mapsto (-1)^{|a||b|}b\hat\otimes a$ для однородных $a,b$.
\end{proof}

\begin{proposition}\label{prop:clifford_superalgebra_of_hyperbolic_plane}
\begin{enumerate}
\item Имеется изоморфизм $C(\mathbb H)\cong\End_F(F\hat\oplus F)$.
\item Если $a,b,c\in F^*$, то
\begin{enumerate}
\item $C(\la ac,bc\ra)\hat\otimes i\left(\quaternion{ac}{bc}{F}\right)\cong C(\la a,b\ra)\hat\otimes i\left(\quaternion{a}{b}{F}\right)$;
\item $C(\lla a,b\rra)\sim i\left(\quaternion{a}{b}{F}\right)$;
\item $C(\lla a,b,c\rra)\sim 1$.
\end{enumerate}
\end{enumerate}
\end{proposition}
\begin{proof}
Заметим, что все участвующие в формулировке супералгебры являются центральными простыми. Поэтому для проверки изоморфизма
достаточно построить гомоморфизм и установить совпадение размерностей; тогда построенный гомоморфизм будет инъективным
в силу простоты и сюръективным из соображений размерности.
\begin{enumerate}
\item Пусть $(e_1,e_2)$~~--- канонический базис гиперболической плоскости
$\mathbb H=\la 1,-1\ra$. Тогда в $C(\mathbb H)$ есть базис
$(1,e_1,e_2,e_1e_2)$ и $|e_1|=|e_2|=1$, $e_1^2=1$, $e_2^2=-1$, $e_1e_2=-e_2e_1$.
Супералгебра $\End_F(F\hat\oplus F)$ обладает базисом $(E_{ij})_{i,j\in\{0,1\}}$ с $|E_{00}|=|E_{11}|=0$
и $|E_{01}|=|E_{10}|=1$. Можно проверить, что искомый изоморфизм индуцируется отображением
$1\mapsto E_{00}+E_{11}$, $e_1\mapsto E_{01}+E_{10}$, $e_2\mapsto E_{01}-E_{10}$.
\item Первое тождество сводится к случаю $ac=1$. Таким образом, достаточно построить гомоморфизм супералгебр
$$
C(\la 1,ab\ra)\hat\otimes i(M_2(F))\stackrel{\sim}\longrightarrow C(\la a,b\ra)\hat\otimes\left(\quaternion{a}{b}{F}\right).
$$
Запишем $M_2(F)$ как алгебру кватернионов с базисом $(1,i,j,ij)$, в котором $i^2=a$, $j^2=1$, $ij=-ji$. Пусть
$(1,i',j',i'j')$~--- канонический базис алгебры $\quaternion{a}{b}{F}$, в котором $i'^2=a$, $j'^2=b$, $i'j'=-j'i'$.
Пусть, наконец, $(e_1,e_2)$ (соответственно $(e'_1,e'_2)$~--- канонический ортогональный базис формы $\la 1,ab\ra$
(соответственно $\la a,b\ra$). Тогда в $C(\la 1,ab\ra)$ (соответственно $C(\la a,b\ra)$) появляется базис
$(1,e_1,e_2,e_1e_2)$ с $e_1^2=1$, $e_2^2=ab$, $e_1e_2=-e_2e_1$ (соответственно 
$(1,e'_1,e'_2,e'_1e'_2)$ с ${e'_1}^2=a$, ${e'_2}^2=b$, $e'_1e'_2=-e'_2e'_1$).
Можно проверить, что отображения
$$
\begin{aligned}
e_1\hat\otimes 1&\mapsto e'_1\hat\otimes i'^{-1}\\
e_2\hat\otimes 1&\mapsto e'_2\hat\otimes i'\\
1\hat\otimes i&\mapsto 1\hat\otimes i'\\
1\hat\otimes j&\mapsto e'_1e'_2\hat\otimes (i'j')^{-1}\\
\end{aligned}
$$
индуцируют требуемый гомоморфизм. Наконец, остальные два пункта следуют из уже доказанного.
\end{enumerate}
\end{proof}

\begin{corollary}\label{corollary:C_homo}
Отображение $q\mapsto C(q)$ индуцирует гомоморфизм
$$
C\colon W(F)/I^3F\to BW(F).
$$
\end{corollary}
\begin{proof}
По теореме~\ref{thm:clifford_superalgebra_is_central_simple} всякая невырожденная форма~$q$
определяет элемент $\la C(q)\ra\in BW(F)$, и по теореме~\ref{thm:tensor_product_of_superalgebras}
имеем $\la C(q\perp q')\ra=\la C(q)\ra + \la C(q')\ra$. Остается применить
предложение~\ref{prop:clifford_superalgebra_of_hyperbolic_plane}.
\end{proof}

Пусть $A=A_0\oplus A_1$~--- центральная простая $F$-супералгебра. Возможны два случая:
\par\noindent$\bullet$ $F$-алгебра $A$ является простой (неградуированной) алгеброй.
Тогда $A_0$ полупроста и ее центр является этальной $F$-алгеброй ранга 2.
В этом случае говорят, что $A$ имеет {\bf четный тип}.
\par\noindent$\bullet$ $A$ не центральна как $F$-алгебра. Тогда $A$ полупроста и ее
центр $Z(A)$ является этальной $F$-алгеброй ранга 2, а $A_0$~--- центральная простая
$F$-алгебра. Кроме того, $A_0$ модуль $A_1$ изоморфен $A_0$. В этом случае говорят,
что $A$ имеет {\bf нечетный тип}.

Определим отображение
$$
\begin{aligned}
\ph\colon BW(F)&\to\mathbb Z/2\times F^*/(F^*)^2\times Br(F)\\
\la A\ra&\mapsto (\varepsilon(A),\delta(A),b(A))\\
\end{aligned}
$$
следующим образом:
\par\noindent$\bullet$ $\varepsilon(A)=\begin{cases}
1,&\text{если $A$ нечетного типа};\\
0,&\text{если $A$ четного типа}.\\
\end{cases}$
\par\noindent$\bullet$ $\delta(A)=\begin{cases}
d(Z(A)),&\text{ если $A$ нечетного типа};\\
d(Z(A_0)),&\text{ если $A$ четного типа},\\
\end{cases}$
\par\noindent
где через $d(E)$ обозначается дискриминант этальной $F$-алгебры $E$.
\par\noindent$\bullet$ $\delta(A)=\begin{cases}
[A_0],&\text{ если $A$ нечетного типа};\\
[A],&\text{ если $A$ четного типа},\\
\end{cases}$

\begin{theorem}\label{thm:BW}
\begin{enumerate}
\item Отображение $\ph$ является биекцией.
\item $\varepsilon$ является гомоморфизмом.
\item Пусть $BW^{(1)}(F)$~--- ядро $\varepsilon$. Ограничение
$\delta$ на $BW^{(1)}(F)$ является гомоморфизмом.
\item Пусть $BW^{(2)}(F)$~--- ядро $\delta$. Тогда
$BW^{(2)}(F)=i(B(F))$ и ограничение $b$ на $BW^{(2)}(F)$
является обратным к $i$.
\item Отображение $\la A\ra\mapsto(\varepsilon(A),\delta(A))$
индуцирует изоморфизм $$BW(F)/BW^{(2)}(F)\to Q(F),$$
где $Q(F)$~--- группа, определенная перед предложением~\ref{prop:group_Q}.
\item Групповой закон, индуцированный на $\mathbb Z/2\times F^*/(F^*)^2\times B(F)$
переносом структуры посредством отображения $\ph$, записывается так:
$$
(m,a,x)+(n,b,y)=(m+n,(-1)^{mn}ab,x+y+((-1)^{m(n+1)}a,(-1)^{(m+1)n}b)),
$$
где через $(u,v)$ обозначается класс алгебры кватернионов $\quaternion{u}{v}{F}$ в $Br(F)$.
\end{enumerate}
\end{theorem}

% 3.05.2010
% когомологии Галуа

\subsection{Когомологии Галуа}

\begin{definition}
Пусть $G$~--- конечная группа. {\bf (Левым) $G$-модулем} называется абелева группа $A$,
снабженная левым действием группы $G$, то есть гомоморфизмом $\ph\colon G\to\Aut(A)$.
Если $A$ записывается аддитивно, то для $(g,a)\in G\times A$ мы будем обозначать
$\ph(g)(a)$ через $ga$. При этом
$$
\begin{aligned}
g(a+b)&=ga+gb,\\
(gh)a&=g(ha).\\
\end{aligned}
$$
\end{definition}
\begin{remarks}
\item Можно определить {\bf правое} действие $G$ на $A$ как {\it анти-гомоморфизм}
из $G$ в $\Aut(A)$; если $A$ записывается аддитивно, то при этом $(g,a)\mapsto ag$.
Задание левого и правого действия $G$ на $A$ эквивалентно: если $\ph$~--- левое действие,
то $g\mapsto \ph(g)^{-1}$~--- правое, и наоборот.
\item Если $A$ записывается мультипликативно, то левое (соответственно правое) действие
удобнее записывать как $(g,a)\mapsto {}^ga$ (соответственно $(g,a)\mapsto a^g$) для избежания
конфликта с записью умножения в $A$.
\item Если $f\colon H\to G$~--- гомоморфизм групп, то на $G$-модуле $A$ появляется структура
$H$-модуля с помощью определения $ha=f(h)a$.
\end{remarks}
\begin{definition}
Пусть $G$~--- конечная группа, $A$~--- левый $G$-модуль с аддитивной записью.
Определим {\bf коцепной комплекс} $G$ со значениями в $A$:
$$
0\to C^0(G,A)\stackrel{d^0}{\longrightarrow}C^1(G,A)\stackrel{d^1}{\longrightarrow}\dots
\stackrel{d^{n-1}}{\longrightarrow} C^n(G,A)\stackrel{d^n}{\longrightarrow}\dots
$$
где $C^n(G,A)$~--- множество всех отображений из $G^n$ в $A$ и
\begin{multline*}
d^nf(g_1,\dots,g_{n+1}=g_1f(g_2,\dots,g_{n+1})\\
+\sum_{j=1}^n(-1)^jf(g_1,\dots,g_jg_{j+1},\dots,g_{n+1})\\
+(-1)^{n+1}f(g_1,\dots,g_n).
\end{multline*}
Элемент $f\in C^n(G,A)$ называется {\bf $n$-коцепью} $G$ со значениями в $A$. Если
$d^nf=0$, $f$ называется {\bf $n$-коциклом}; подгруппа $n$-коциклов обозначается
через $Z^n(G,A)$. Если $f\in\Image d^{n-1}$, $f$ называется {\bf $n$-кограницей};
подгруппа $n$-кограниц обозначается через $B^n(G,A)$.
\end{definition}

\begin{examples}
\item 0-коцикл~--- это элемент $A^G:=\{a\in A\mid ga=a\;\forall g\in G\}$.
\item 1-коцикл~--- это отображение $f\colon G\to A$ такое, что $f(gh)=f(g)+gf(h)$.
Такое отображение называется {\bf скрещенным гомоморфизмом}.
\item 2-коцикл~--- это отображение $f\colon G\times G\to A$ такое, что
$$
f(g,h)+f(gh,k)=fg(h,k)+f(g,hk).
$$
Такое $f$ называется {\bf системой факторов}.
\end{examples}

\begin{lemma}
Если $n>0$, то $d^{n+1}d^n=0$; иными словами, $B^n(G,A)\subset Z^n(G,A)$.
\end{lemma}
\begin{proof}
Простое вычисление.
\end{proof}

\begin{definition}
{\bf $n$-ой группой когомологий $G$ со значениями в $A$} называется фактор-группа
$H^n(G,A)=Z^n(G,A)/B^n(G,A)$.
\end{definition}

\begin{proposition}
\begin{enumerate}
\item Функтор $(G,A)\mapsto H^n(G,A)$ ковариантен по $A$ и контравариантен по $G$.
Если $f\colon H\to G$~--- гомоморфизм групп, обозначим через $f^*$ индуцированный
гомоморфизм групп когомологий.
\item Пусть $0\to A'\to A\to A''\to 0$~--- короткая точная последовательность $G$-модулей.
Тогда существует длинная точная последовательность
\begin{multline*}
0\to H^0(G,A')\to H^0(G,A)\to H^0(G,A'')\to H^1(G,A')\to\dots\\
\to H^n(G,A')\to H^n(G,A)\to H^n(G,A'')\to H^{n+1}(G,A)\to\dots
\end{multline*}
\item Если действие $G$ на $A$ тривиально, то существует канонический изоморфизм
между $H^1(G,A)$ и $\Hom(G,A)$.
\end{enumerate}
\end{proposition}

\begin{definition}
\begin{enumerate}
\item Пусть $H\leq G$. Отображение $H^*(G,A)\to H^*(H,A)$, индуцированное вложением $H$ в $G$,
называется {\bf морфизмом ограничения} и обозначается через $\Res^H_G$.
\item Пусть $H\to G$~--- сюръективный гомоморфизм групп. Индуцированное им отображение $H^*(G,A)\to H^*(H,A)$
называется {\bf морфизмом инфляции} и обозначается через $\Inf^H_G$.
\end{enumerate}
\end{definition}

\begin{theorem}
Пусть $G$~--- конечная группа, $H\leq G$.
\begin{enumerate}
\item Существует единственный набор гомоморфизмов
$$
\Cor^G_H\colon H^n(H,A)\to H^n(G,A)
$$
(для любого
$G$-модуля $A$ и для любого $n\geq 0$), естественных по $A$, согласованных с длинными точными последовательностями,
ассоциированными с короткими точными последовательностями $G$-модулей, и таких, что в степени 0
$$
\Cor^G_H(a)=\sum_{g\in G/H}ga
$$
для всех $a\in H^0(H,A)=A^H$.
\item Если $m=(G:H)$~--- индекс $H$ в $G$, то $\Cor^G_H\circ\Res^H_G=m$.
\end{enumerate}
\end{theorem}

\begin{definition}
Пусть $A$, $B$~--- $G$-модули. {\bf Тензорным произведением} $A$ и $B$ называется
абелева группа $A\otimes B$, снабженная диагональным действием $G$: $g(a\times b)=ga\times gb$.
\end{definition}

\begin{theorem}
Пусть $A,B$~--- два $G$-модуля. Существуют билинейные гомоморфизмы
$$
\begin{aligned}
H^p(G,A)\times H^q(G,B)&\to H^{p+q}(G,A\otimes B),\quad p,q\geq 0\\
(x,y)&\mapsto x\cdot y,
\end{aligned}
$$
естественные по $A$ и $B$. Они обладают следующими свойствами:
\begin{enumerate}
\item {\bf Ассоциативность}: если $C$~--- еще один $G$-модуль и
$x\in H^p(G,A)$, $y\in H^q(G,B)$, $z\in H^r(G,C)$, то
$(x\cdot y)\cdot z=x\cdot(y\cdot z)$ с учетом изоморфизма
$(A\otimes B)\otimes C\stackrel{\sim}{\to}A\otimes(B\otimes C)$,
при котором $(a\otimes b)\otimes c\mapsto a\otimes(b\otimes c)$.
\item {\bf Коммутативность}: если $x\in H^p(G,A)$, $y\in H^q(G,B)$,
то $x\cdot y=(-1)^{pq}y\cdot x$ с учетом изоморфизма
$A\otimes B\stackrel{\sim}{\to}B\otimes A$, при котором $a\otimes b\mapsto b\otimes a$.
\item {\bf Контравариантность по $G$}: если $f\colon H\to G$~--- гомоморфизм групп, то
$f^*(x\cdot y)=f^*x\cdot f^*y$ для всех $x\in H^p(G,A)$, $y\in H^q(G,B)$.
\item {\bf Формула проекции}: если $H$~--- подгруппа $G$, то
$\Cor^G_H(x\cdot\Res^H_G y)=(\Cor^G_Hx)\cdot y$ для всех $x\in H^p(H,A)$, $y\in H^q(G,B)$.
\end{enumerate}
Таким образом введенное произведение на когомологиях называется {\bf чашечным произведением}.
\end{theorem}

Чашечное произведение можно определить как билинейное отображение на коцепях:
если $f\in C^m(G,A)$, $f'\in C^n(G,B)$, то
$$
(f\cdot f')(g_1,\dots,g_{m+n})=f(g_1,\dots,g_m)\otimes g_1\dots g_m f'(g_{m+1},\dots,g_{m+n}).
$$

\begin{definition}
Топологическая группа $G$ называется {\bf проконечной}, если она удовлетворяет одному из следующих эквивалентных условий:
\begin{enumerate}
\item $G$ является проективным пределом конечных групп;
\item $G$ отделима и всякая открытая подгруппа $G$ имеет конечный индекс.
\end{enumerate}
\end{definition}

\begin{definition}
Пусть $G$~--- проконечная группа. {\bf Топологический $G$-модуль}~--- это абелева группа $A$, снабженная
левым действием группы $G$, такая, что $A=\bigcup_HA^H$, где объединение берется по всем открытым
(следовательно, замкнутым и конечного индекса) подгруппам $G$.
Если $G$~--- проконечная группа, $A$~--- топологический $G$-модуль, определим
{\bf группы когомологий $G$ с коэффициентами в $A$} формулой
$$
H^n(G,A)=\varinjlim H^n(G/H,A^H),
$$
где $H$ пробегает все различные открытые подгруппы в $G$, а морфизмы, участвующие в
определении предела~--- морфизмы инфляции.
\end{definition}

Можно доказать, что когомологии проконечных групп обладают теми же основными свойствами, что и когомологии
конечных групп.
\begin{definition}
Пусть $F_s$~--- сепарабельное замыкание $F$. Группа $F$-автоморфизмов поля $F_s$ обладает
структурой проконечной группы:
$$
G_F=\varprojlim Gal(E/F),
$$
где $E/F$ пробегает все конечные подрасширения Галуа в $F_s$. Группа $G_F$ называется
{\bf абсолютной группой Галуа} поля $F$; ее когомологии называются {\bf когомологиями Галуа};
обычно мы пишем $H^*(F,A)$ вместо $H^*(G_F,A)$.
\end{definition}

\begin{lemma}\label{lemma:automorphisms_are_linearly_independent}
Автоморфизмы поля $E$ являются линейно независимыми над $E$ отображениями.
\end{lemma}
\begin{proof}
Допустим, что $\sum a_\ph\ph=0$, где $\ph$~--- автоморфизмы $E$, $a_\ph\in E$.
Можно предположить, что множество ненулевых коэффициентов $a_\ph$ имеет минимально возможную мощность.
В этом множестве хотя бы два элемента; возьмем $\ph_1\neq\ph_2$ такие, что $a_{\ph_1},a_{\ph_2}\neq 0$.
Найдется $x\in E$ такой, что $\ph_1(x)\neq\ph_2(x)$. Тогда для любого $y\in E$ имеем
$$
\sum a_\ph\ph(y)=0,
$$
откуда
$$
0=\sum a_\ph\ph(xy)-\ph_1(x)\sum a_\ph\ph(y)=\sum a_\ph(\ph(x)-\ph_1(x))\ph(y),
$$
поэтому $\sum a_\ph(\ph(x)-\ph_1(x))\ph=0$~--- новая линейная зависимость, в которой меньше ненулевых
слагаемых, чем в исходной: противоречие.
\end{proof}

\begin{theorem}{Гильберта 90}
Пусть $E/F$~--- расширение Галуа с группой $G$. Тогда $H^1(G,E^*)=0$.
\end{theorem}
\begin{proof}
Пусть $(a_g)_{g\in G}$~--- 1-коцикл $G$ со значениями в $E^*$.
По лемме~\ref{lemma:automorphisms_are_linearly_independent} найдется $x\in E$
такой, что $a=\sum_{g\in G}a_g{}^gx\neq 0$. Поэтому для любого $h\in G$
$$
{}^ha=\sum_{g\in G}{}^ha_g{}^{hg}x=\sum_{g\in G}a_h^{-1}a_{hg}{}^{hg}x= a_h^{-1}\sum_{g\in G}a_g{}^gx=a_h^{-1}a,
$$
что и означает, что $(a_g)$ является 1-кограницей.
\end{proof}
\begin{corollary}
$H^1(F,F^*_s)=0$.
\end{corollary}

% 10.05.2010

\begin{definition}
Пусть $E/F$~--- расширение Галуа с группой $G$. Пусть $G$ действует слева на $E$,
$c$~--- 2-коцикл $G$ со значениями в $E^*$. {\bf Скрещенным произведением},
соответствующим $c$, называется следующая $F$-алгебра $E\times_sG$:
\par\noindent$\bullet$ {\it Аддитивная структура:} $E\times_cG$~--- векторное пространство над $E$
с базисом $G$.
\par\noindent$\bullet$ {\it Мультипликативная структура:} умножение $F$-билинейно и если
$x,y\in E$, $g,h\in G$, то
$$
(x\cdot g)(y\cdot h)=x{}^gyc_{g,h}\cdot gh.
$$
\end{definition}

\begin{theorem}
\begin{enumerate}
\item Таким образом определенная алгебра $E\times_cG$ является ассоциативной и центральной простой над $F$
степени $n=[E:F]$; $E$~--- максимальная коммутативная подалгебра $E\times_cG$.
\item Всякая центральная простая $F$-алгебра $A$, содержащая $E$ как максимальное коммутативное под-тело,
имеет вид $E\times_cG$.
\item Пусть $c,c'$~--- коциклы $G$ со значениями в $E^*$. $E\times_cG\cong E\times_{c'}G$ тогда и только тогда,
когда $c$ и $c'$ когомологичны (то есть $c/c'\in B^2(G,E^*)$.
\end{enumerate}
\end{theorem}

\begin{corollary}\label{corollary:secong_cohomology_as_brauer_group}
Существует канонический изоморфизм
$$
u_{E/F}\colon H^2(G,E^*)\stackrel{\sim}{\to} Br(E/F),
$$
где $Br(E,F)=\Ker(Br(F)\to B(E))$.
\end{corollary}

\begin{example}\label{example:second_cohomology_of_quadratic_extension}
$E=F(\sqrt{a})$, где $a\notin F^*/(F^*)^2$. Тогда $G=\{1,g\}$. Будем обозначать
действие $g$ через $x\mapsto\overline{x}$. 2-коцикл $G$ с коэффициентами в $E^*$
задается четырьмя элементами $c_{1,1}$, $c_{1,g}$, $c_{g,1}$ и $c_{g,g}$.
Применяя соотношение коцикла к тройкам $(1,1,g)$, $(g,1,1)$ и $(g,g,g)$, получаем
$$
\begin{aligned}
c_{1,1}&=c_{1,g}\\
c_{g,1}&=\overline{c_{1,1}}\\
c_{g,g}c_{1,g}&=\overline{c_{g,g}}c_{g,1}.
\end{aligned}
$$
После деления на кограницу можно считать, что $c_{1,1}=1$ (такой 2-коцикл называют
{\bf нормализованным}). Тогда $c_{1,g}=c_{g,1}=1$ и $c_{g,g}=b\in F^*$.
Пусть $\alpha\in E^*$ и $\alpha^2=a$; положим $\beta=\alpha\cdot g\in A$.
Тогда $\alpha\beta=-\beta\alpha$ и $\beta^2=-ab$; значит, мы получили алгебру
кватернионов $\quaternion{a}{-ab}{F}=\quaternion{a}{b}{F}$.
\end{example}
\begin{theorem}
Изоморфизмы $u_{E/F}$ из следствия~\ref{corollary:secong_cohomology_as_brauer_group}
склеиваются в изоморфизм
$$
u_F\colon H^2(F,F^*_s)\stackrel{\sim}{\to}Br(F).
$$
\end{theorem}

Пусть $n$~--- натуральное число, взаимно простое с характеристикой $F$; тогда возведение в степень $n$
сюръективно на $F_s^*$. Рассмотрим {\bf точную последовательность Куммера} $G_F$-модулей
$$
1\to\mu_n\to F_s^*\stackrel{n}{\longrightarrow}F_s^*\to 1,
$$
где $\mu_n$~--- группа корней $n$-ой степени из 1 в $F_s$. Ей соответствует длинная
точная последовательностей когомологий Галуа:
\begin{multline*}
1\to H^0(F,\mu_n)\to H^0(F,F_s^*)\stackrel{n}{\to}H^0(F,F_s^*)\\
\stackrel{\delta}{\to}H^1(F,\mu_n)\to H^1(F,F_s^*)\stackrel{n}{\to}H^1(F,F_s^*)\\
\to H^2(F,\mu_n)\to H^2(F,F_s^*)\stackrel{n}{\to}H^2(F,F_s^*).
\end{multline*}
Заметим, что $H^0(F,F_s^*)=F^*$ и $H^1(F,F_s^*)=0$ по теореме Гильберта 90.
Мы получили следующую теорему:
\begin{theorem}[теория Куммера]\label{thm:Kummer}
Эта точная последовательность приводит к изоморфизмам
$$
\begin{aligned}
F^*/(F^*)^n&\stackrel{\sim}{\to}H^1(F,\mu_n)\\
H^2(F,\mu_n)&\stackrel{\sim}{\to}{}_nBr(F).
\end{aligned}
$$
\end{theorem}

Первый изоморфизм мы будем обозначать через $a\mapsto (a)$.
В случае $n=2$ группа $\mu_2$ является тривиальным $G_F$-модулем,
изоморфным группе $\mathbb Z/2$. Получаем изоморфизмы
$$
\begin{aligned}
F^*/(F^*)^2&\stackrel{\sim}{\to}H^1(F,\mathbb Z/2)\\
H^2(F,\mathbb Z/2)&\stackrel{\sim}{\to}{}_2Br(F).
\end{aligned}
$$
В частности, если $a,b\in F^*$, то класс $(a,b)$ алгебры кватернионов,
определенной элементами $a$ и $b$ является элементом $H^2(F,\mathbb Z/2)$.

\begin{proposition}\label{prop:cup_product}
$(a,b)=(a)\cdot(b)$.
\end{proposition}
\begin{proof}
По определению чашечное произведение задается формулой
$(g,h)\mapsto (a)(g)\cdot(b)(h)$, и его образ в $H^2(F,F_s^*)$ представляется 2-коциклом
$b_{g,h}=(-1)^{(a)(g)\cdot (b)(g)}$. В примере~\ref{example:second_cohomology_of_quadratic_extension}
мы видели, что класс алгебры $\quaternion{a}{b}{F}$ в $H^2(F,F_s^*)$ представляется
2-коциклом
$$
c_{g,h}=\begin{cases}
1,&\text{если $(a)(g)=0$ или $(a)(h)=0$};\\
-ab,&\text{если $(a)(g)=(a)(h)=1$}.\\
\end{cases}
$$
Осталось проверить, что $b_{g,h}$ и $c_{g,h}$ когомологичны.
Возьмем $\alpha,\beta\in F_s^*$ такие, что $\alpha^2=a,\beta^2=b$.
Легко видеть, что $b_{g,h}^{-1}c_{g,h}=f(gh)^{-1}f(g){}^gf(h)$, где
$$
f(g)=\begin{cases}
1,&\text{если $(a)(g)=0$};\\
\alpha\beta,&\text{если $(a)(g)=1$, $(b)(g)=0$};\\
-\alpha\beta,&\text{если $(a)(g)=(b)(g)=1$}.\\
\end{cases}
$$
\end{proof}

\subsection{Теорема Меркурьева}

\begin{proposition}\label{proposition:Merkurjev}
Пусть $q$~--- квадратичная форма над $F$. Тогда
$\varepsilon(C(q))=\overline{\dim q}$ и $\delta(C(q))=d(q)$
(см. обозначения перед теоремой~\ref{thm:BW}.
\end{proposition}
\begin{proof}
Рассмотрим гомоморфизмы $(\overline{\dim},d)$ и $(\varepsilon,\delta)\circ C$
из $W(F)/I^3F$ в $Q(F)$. Для доказательства их совпадения достаточно проверить это
на порождающих $W(F)$, скажем, на классах одномерных форм $\la a\ra$, $a\in F^*$,
что очевидно.
\end{proof}

\begin{lemma}
Пусть $q\in I^2F$. Тогда $C_0(q)\cong A\times A$ и $C(q)\cong M_2(A)$ для некоторой
центральной простой алгебры $A$.
\end{lemma}
\begin{proof}
По предложению~\ref{proposition:Merkurjev} алгебра $C(q)$ имеет четный тип и $\delta(C(q))=1$.
По теореме~\ref{thm:BW} получаем, что $C(q)$ подобна алгебре $i(A_0)$ для некоторой центральной
простой алгебры $A_0$. Другими словами, существует векторное суперпространство $V=V_0\oplus V_1$
такое, что $C(q)\cong i(A_0)\hat\otimes_F\End_F(V)$. Размерности $C_0(q)$ и $C_1(q)$ совпадают,
поэтому $\dim V_0=\dim V_1$. Отождествляя $V_1$ с $V_0$, получаем, что
$$
C(q)\cong i(A_0)\hat\otimes_F\End_F(V)\cong i(A_0\otimes_F\End_F(V_0))\hat\otimes_F M_2(F),
$$
и получаем нужное утверждение для $A=A_0\otimes_F\End_F(V_0)$.
\end{proof}

%\noindent
$$
\begin{tabular}{|c|c|c|c|c|c|c|}
\hline
$\dim q$ & $d(q)$ & $Z(C(q))$ & $C(q)$ & $Z(C_0(q))$ & $C_0(q)$ & $\deg(q)$\\
\hline
\hline
\multirow{2}{*}{нечетно} & $\notin (F^*)^2$ & $F(\sqrt{d})$ & простая &
  \multirow{2}{*}{$F$} & центральная & \multirow{2}{*}{0}\\
\cline{2-4}
& $\in(F^*)^2$ & $F\times F$ & $C_0(q)\times C_0(q)$ & & простая&\\
\hline
\multirow{2}{*}{четно} & $\notin(F^*)^2$ & \multirow{2}{*}{$F$} & цпа & $F(\sqrt{d})$ & простая & 1\\
\cline{2-2}\cline{4-7}
& $\in(F^*)^2$ & & $M_2(A)$ & $F\times F$ & $A\times A$, $A$ цпа & $\geq 2$\\
\hline
\end{tabular}
$$

\begin{definition}
Для квадратичной формы $q$ обозначим через $c(q)$ элемент $b(C(q))\in Br(F)$~---
{\bf инвариант Клиффорда} $q$. Таким образом,
$$
c(q)=\begin{cases}
[C_0(q)],&\text{ если $A$ нечетного типа};\\
[C(q)],&\text{ если $A$ четного типа}.\\
\end{cases}
$$
\end{definition}

\begin{proposition}\label{prop:c_homo}
Пусть $q,q'$~--- две квадратичные формы. Тогда
$$
c(q\perp q')=c(q)+c(q')+((-1)^{m(n+1)}d(q),(-1)^{(m+1)n}d(q')),
$$
где $m=\dim q$, $n=\dim q'$.
\end{proposition}
\begin{proof}
Тривиально следует из теоремы~\ref{thm:BW}.
\end{proof}

\begin{proposition}\label{prop:c_otimes}
\begin{enumerate}
\item Пусть $\ph,\psi\in IF$. Тогда $c(\ph\otimes\psi)=(d(\ph),d(\psi))$.
\item Пусть $q$~--- квадратичная форма, $a\in F^*$. Тогда
$$
c(aq)=\begin{cases}
c(q)+(a,d(q)),&\text{ если $\dim q$ четна};\\
[c(q)],&\text{ если $\dim q$ нечетна}.\\
\end{cases}
$$
\end{enumerate}
\end{proposition}


\begin{lemma}
\begin{enumerate}
\item Пусть $a,b\in F^*$. Тогда $\lla a,b\rra$ гиперболична $\Longleftrightarrow$
$c(\lla a,b\rra)=0$.
\item Пусть $a,b,c,d\in F^*$. Тогда $\lla a,b\rra\cong\lla c,d\rra$ $\Longleftrightarrow$ $(a,b)=(c,d)$.
\item Пусть $\sigma,\tau\in GP_2(F)$. Тогда  $\sigma$ пропорциональна $\tau$ $\Longleftrightarrow$ $c(\sigma)=c(\tau)$.
\end{enumerate}
\end{lemma}
\begin{proof}
(1)~--- очевидно.
Для доказательства (2) предположим сначала, что $b=d$. Тогда $(ac,b)=0$, поэтому $\lla ac,b\rra\sim 0$ по пункту $(1)$.
Значит, $\lla a,b\rra\perp -\lla c,d\rra\sim c\lla ac,b\rra\sim 0$, что и требовалось. В общем случае применим
лемму Альберта~\ref{lemma:chain_albert}: найдется $e$ такое, что $(a,b)=(a,e)=(c,e)=(c,d)$, поэтому
$\lla a,b\rra\cong\lla a,e\rra\cong\lla c,e\rra\cong\lla c,d\rra$.
В (3) пусть $\sigma=a\sigma_0$, $\tau=b\tau_0$, где $\sigma_0,\tau_0\in P_2(F)$. Тогда $c(\sigma)=c(\sigma_0)$,
$c(\tau)=c(\tau_0)$, и все следует из пункта $(2)$.
\end{proof}

Из предложения~\ref{prop:c_homo} следует, что ограничение инварианта Клиффорда $c$ на $I^2F$ является
гомоморфизмом, принимающим значения в подгруппе 2-кручения ${}_2Br(F)$ группы Брауэра $Br(F)$.
\begin{theorem}[Меркурьев]\label{thm:Merkurjev}
Гомоморфизм
$$
c\colon I^2F/I^3F\to {}_2Br(F),
$$
индуцированный инвариантом Клиффорда, является изоморфизмом.
\end{theorem}

Обозначим через $BW_2(F)$ множество элементов $x\in BW(F)$ таких, что $b(x)\in{}_2Br(F)$. Легко проверить,
что $BW_2(F)$ является подгруппой в $BW(F)$, содержащей $i({}_2Br(F))$ (но она не совпадает с
подгруппой 2-кручения $BW(F)$!). Из теоремы Меркурьева нетрудно вывести следующее утверждение.
\begin{corollary}
Гомоморфизм $C$ (см. следствие~\ref{corollary:C_homo}) индуцирует изоморфизм
$$
C\colon W(F)/I^3F\to BW_2(F).
$$
\end{corollary}

\subsection{Высшие инварианты}

Мы будем обозначать через $H^nF$ группы когомологий Галуа
$H^n(F,\mathbb Z/2)$. Мы знаем, что инварианты
$\overline{\dim}$, $d$ и $c$ можно рассматривать как
инварианты
$$
e^n\colon W(F)\to H^nF
$$
для $n=0,1,2$ (см. теорему~\ref{thm:Kummer} и замечание после нее).

\begin{theorem}
Для $n\leq 2$ инвариант $e^n$ индуцирует изоморфизм
$$
\overline{e}^n\colon I^nF/I^{n+1}F\to H^nF,
$$
причем
$$
e^{p+q}(xy)=e^p(x)\cdot e^q(y)
$$
для $p+q\leq 2$, $x\in I^pF/I^{p+1}F\times I^qF/I^{q+1}F$.
\end{theorem}
\begin{proof}
Первое утверждение являеется переформулировкой уже известных
теорем~\ref{thm:discriminant_iso} и~\ref{thm:Merkurjev}.
Второе нужно проверить только для $p=q=1$, а это следует
из предложений~\ref{prop:c_otimes} и~\ref{prop:cup_product}.
\end{proof}

Напомним, что {\bf $K$-теорией Милнора} называется
градуированное кольцо $K_*^M(F)$, заданное образующими
$\{a\}$, $a\in F^*$ и соотношениями
$\{ab\}=\{a\}+\{b\}$ ($a,b\in F^*$),
$\{a\}\cdot\{1-a\}=0$ ($a\in F^*\setminus\{1\}$).
Иными словами, $K_*^M(F)$ является фактором тензорной алгебры
$\mathbb Z$-модуля $F^*$ по двустороннему идеалу, порожденному
элементами $a\otimes(1-a)$ для $a\neq 1$. Легко видеть, что
$K_0(F)=\mathbb Z$, $K_1(F)=F^*$. Будем обозначать произведение
$\{a_1\}\cdot\dots\{a_n\}\in K_n^M(F)$ через $\{a_1,\dots,a_n\}$.

\begin{lemma}\label{lemma:simple_k_theory}
В кольце $K_*^M(F)$ выполнены соотношения $\{a,a\}=\{a,-1\}$
и $\{a,b\}=-\{b,a\}$ для всех $a,b\in F^*$.
\end{lemma}
\begin{proof}
Поскольку $1$ является нейтральным элементом абелевой группы $F^*$,
имеем $\{1,1\}=\{1,-1\}=0$. Пусть теперь $a\neq 0,1$. Тогда
$\{a,1-a\}=0=\{a^{-1},1-a^{-1}\}$. Из билинейности следует, что
$\{a^{-1},1-a^{-1}\}=-\{a,1-a^{-1}\}=-\{a,\frac{1-a}{-a}\}=
-\{a,1-a\}+\{a,-a\}$, поэтому $\{a,-a\}=0$.
Но $-a=\frac{a}{-1}$, поэтому $\{a,a\}=\{a,-1\}$.
Для доказательства второго соотношения заметим, что
$\{ab,ab\}=\{a,a\}+\{b,b\}+\{a,b\}+\{b,a\}=
\{a,-1\}+\{b,-1\}+\{a,b\}+\{b,a\}=
\{ab,-1\}+\{a,b\}+\{b,a\}=
\{ab,ab\}+\{a,b\}+\{b,a\}$ по уже доказанному;
отсюда $\{a,b\}+\{b,a\}=0$.
\end{proof}

\begin{examples}
\begin{enumerate}
\item Пусть $F=\mathbb F_q$. Известно, что мультипликативная группа конечного поля является
циклической, поэтому $K_1^M(\mathbb F_q)\cong\mathbb Z/(q-1)$. Умножение в $K_*^M(\mathbb F_q)$
дает нам сюръективный гомоморфизм
$$
\mathbb F_q^*\otimes_{\mathbb Z}\mathbb F_q^*\to K_2^M(\mathbb F_q).
$$
Пусть $a$~--- образующая группы $\mathbb F_q^*$; тогда $K_2^M(\mathbb F_q)$ также является
циклической с образующей $\{a,a\}$. Но по лемме~\ref{lemma:simple_k_theory} $\{a,a\}=\{a,-1\}$,
стало быть, эта образующая имеет порядок 1 или 2. Значит, во всяком случае,
$$
K_2^M(\mathbb F_q)\cong K_2^M(\mathbb F_q)/2.
$$
Заметим, что уравнение $x^2+y^2=a$ имеет нетривиальное решение над $\mathbb F_q$. Можно считать,
что $x\neq 0$, тогда $1+y^2/x^2=a/x^2$ и
$$
\begin{aligned}
\{a,-1\}&=\{a/x^2,-1\}\\
&=\{a/x^2,-1\}+\{a/x^2,1-a/x^2\}\\
&=\{a/x^2,a/x^2-1\}\\
&=\{a/x^2,y^2/x^2\}\\
&=\{a/x^2,1\}\\
&=0
\end{aligned}
$$
в группе $K_2^M(\mathbb F_q)/2$, значит, и в $K_2^M(\mathbb F_q)$. Поэтому $K_2^M(\mathbb F_q)=0$
и, следовательно, $K_n^M(\mathbb F_q)=0$ для всех $n\geq 2$.
\item Пусть $F=\mathbb R$; тогда $\mathbb R^*/2\cong\mathbb Z/2$~--- циклическая группа порядка 2
с образующей $-1$. Значит, $T_{\mathbb Z/2}(\mathbb R^*/2)=(\mathbb Z/2)[t]$, где $t=\{-1\}$.
С другой стороны, из двух элементов $a,1-a$ хотя бы один является положительным, поэтому из него
извлекается квадратный корень в $\mathbb R$. Это означает, что элемент
$a\otimes(1-a)$ является 2-делимым в $(\mathbb R^*/2)^{\otimes 2}$.
Поэтому $K_*^M(\mathbb R)/2\cong T_{\mathbb Z/2}(\mathbb R^*/2)\cong(\mathbb Z/2)[t]$, где $t=\{-1\}$.
\end{enumerate}
\end{examples}

\begin{proposition}
Пусть $n\geq 0$. Существуют гомоморфизмы
$$
\begin{aligned}
a^n\colon & K_n^M(F)/2\to I^nF/I^{n+1}F,\\
b^n\colon & K_n^M(F)/2\to H^nF
\end{aligned}
$$
такие, что
$$
\begin{aligned}
a^n(\{a_1,\dots,a_n\}&=\lla a_1,\dots,a_n\rra\\
b^n(\{a,1,\dots,a_n\}&=(a_1)\cdot\dots\cdot(a_n).
\end{aligned}
$$
Кроме того, гомоморфизм $a^n$ является сюръективным.
Произведение $(a_1)\cdot\dots\cdot(a_n)$ в $H^*F$ мы будем обозначать через $(a_1,\dots,a_n)$.
\end{proposition}
\begin{proof}
Докажем, что $a^n$ и $b^n$~--- корректно определенные отображения.
Сопоставление
$(a_1,\dots,a_n)\mapsto\lla a_1,\dots,a_n\rra$ является полилинейным
по лемме~\ref{lemma:pfister_identity}. Форма $2=\la 1,1\ra$ лежит
в $IF$, поэтому $2I^nF/I^{n+1}F=0$. Форма $\lla a,1-a\rra$ изотропна,
поэтому она эквивалентна 0.
Отображение
$(a_1,\dots,a_n)\mapsto(a_1)\cdot\dots\cdot(a_n)$ является полилинейным
в силу определения; кроме того, очевидно, что $2H^nF=0$. Наконец
из леммы~\ref{lemma:relations_quaternion} и предложения~\ref{prop:cup_product}
следует, что $(a)\cdot(1-a)=0$.
Сюръективность $a^n$ следует из того, что $I^nF/I^{n+1}F$ порождается
классами $n$-форм Пфистера
\end{proof}

\begin{proposition}[Милнор]
Существует отображение
$$
w\colon \tilde W(F)\to K_*^M(F)/2
$$
такое, что $w(q\perp q')=w(q)w(q')$ для $q,q'\in\tilde W(F)$ и $w(\la a\ra)=1+\{a\}$.
Кроме того,
$$
w((\la 1\ra-\la a_1\ra)\otimes\dots\otimes(\la 1\ra-\la a_1\ra))=1+\{-1\}^{2^{n-1}-n}\{a_1,\dots,a_n\}.
$$
Для $q\in\tilde W(F)$ запишем $w(q)=\sum_{n\geq 0}w_n(q)$, где $w_n(q)\in K_n^M(F)/2$. Классы
$w_n(q)$ называются {\bf классами Штифеля--Уитни} формы $q$.
\end{proposition}

\begin{proof}
В силу предложения~\ref{theorem:witt_ring_generators_and_relations}, для доказательства существования $w$
достаточно проверить, что $w(\la a,b\ra)=w(\la a+b,ab(a+b)\ra)$ для всех $a,b\in F^*$, $a+b\neq 0$.
По определению $w(\la a,b\ra)=(1+\{a\})(1+\{b\})=1+\{ab\}+\{a,b\}$ и
$w(\la a+b,ab(a+b)\ra)=(1+\{a+b\})(1+\{ab(a+b)\ra)=1+\{ab\}+\{a+b,ab(a+b)\}$.
Осталось заметить, что $\{a+b,ab(a+b)\}=\{a+b,-ab\}=\{a,-ab\}+\{1+b/a,-ab\}=\{a,b\}+\{a+b/a,-b/a\}=\{a,b\}$.

Положим $\lla a_1,\dots,a_n\rra^{\tilde{}}:=(\la 1\ra-\la a_1\ra)\otimes\dots\otimes(\la 1\ra-\la a_n\ra)$.
Для доказательства последней формулы проведем индукцию по $n$. Очевидно, что
$$
\lla a_1,\dots,a_n\rra^{\tilde{}}=\lla a_1,\dots,a_{n-1}\rra^{\tilde{}}-a_n\lla a_1,\dots,a_{n-1}\rra^{\tilde{}}.
$$
Значит,
$$
w(\lla a_1,\dots,a_n\rra^{\tilde{}})=w(\lla a_1,\dots,a_{n-1}\rra^{\tilde{}})w(a_n\lla a_1,\dots,a_{n-1}\rra^{\tilde{}})^{-1}.
$$
Обозначим $X_n=\{-1\}^{2^{n-1}-n}\{a_1,\dots,a_n\}$.
Нетрудно видеть, что если $q\in \tilde W(F)$~--- произвольная форма размерности $m$, $a\in F^*$ и
$$
w(q)=\sum_{i\geq 0} w_i(q),\quad w_i(q)\in K_i^M(F)/2.
$$
то
$$
w(aq)=\sum_{i\geq 0}(1+\{a\})^{m-i}w_i(q).
$$
Поэтому
$$
w(a_n\lla a_1,\dots,a_{n-1}\rra^{\tilde{}})=1+(1+\{a_n\})^{-2^{n-2}}X_{n-1}.
$$
Значит,
$$
\begin{aligned}
w(\lla a_1,\dots,a_n\rra^{\tilde{}})&=(1+X_{n-1})(1+(1+\{a_n\})^{-2^{n-2}}X_{n-1})^{-1}\\
&=(1+X_{n-1})(1+\{a_n\})^{2^{n-2}}\left((1+\{a_n\})^{2^{n-2}}+X_{n-1}\right)^{-1}\\
&=(1+X_{n-1}+\{a_n\}^{2^{n-2}}+X_{n-1}\{a_n\}^{2^{n-2}})(1+\{a_n\}^{2^{n-2}}+X_{n-1})^{-1}.
\end{aligned}
$$
Заметим, что $\{a_n\}^{2^{n-2}}=\{-1\}^{2^{n-2}-1}\{a_n\}$ (по лемме~\ref{lemma:simple_k_theory}),
откуда
$X_{n-1}\{a_n\}^{2^{n-2}}=X_n$ и
$$
w(\lla a_1,\dots,a_n\rra^{\tilde{}})=1+X_n(1+\{a_n\}^{2^{n-2}}+X_{n-1})^{-1}.
$$
Обозначим $A:=\{a_n\}^{2^{n-2}}$, $B:=X_{n-1}$. Тогда
$$
X_n(\{a_n\}^{2^{n-2}}+X_{n-1})=A^2B+AB^2=\{-1\}^{2^{n-2}}AB+\{-1\}^{2^{n-2}}AB=0,
$$
откуда $w(\lla a_1,\dots,a_n\rra^{\tilde{}})=1+X_n$.
\end{proof}

\end{document}


